\begin{savequote}
Ein Mann ging von Jerusalem nach Jericho hinab und wurde von Räubern überfallen. 
Sie plünderten ihn aus und schlugen ihn nieder; dann gingen sie weg und ließen ihn 
halbtot liegen. Zufällig kam ein Priester denselben Weg
herab; er sah ihn und ging weiter. Auch ein Levit kam zu
der  Stelle; er sah ihn und ging weiter. Dann kam ein Mann
aus  Samarien, der auf der Reise war. Als er ihn sah, hatte
er  Mitleid, ging zu ihm hin, goss Öl und Wein auf seine
Wunden  und verband sie. 

Dann hob er ihn auf sein Reittier,
brachte  ihn zu einer Herberge und sorgte für ihn. Am
anderen Morgen  holte er zwei Denare hervor, gab sie dem
Wirt und sagte:  Sorge für ihn, und wenn du mehr für ihn
brauchst, werde  ich es dir bezahlen, wenn ich wiederkomme.
Was meinst du:  Wer von diesen dreien hat sich als der
Nächste dessen erwiesen,  der von den Räubern überfallen
wurde?  

Der Gesetzeslehrer antwortete: Der, der barmherzig
an  ihm gehandelt hat. Da sagte Jesus zu ihm: Dann geh und handle genauso!
\qauthor{Lk 10,25-37}
\end{savequote}

\chapter{Krankheiten und Medizinische Notfälle --- \gAasRsN{RS VI}}



\minitoc

\author{Sebastian Gabriel, Michael Withofner}

\begin{verbatim}

Changelog:

2009-03-14: In LaTeX überführt

2009-02-25: Start Changelog. Bis inkl. bakt. Meningitis mit Korr. von
Koch verglichen und korr.
\end{verbatim}


\section{Einführung in die Medizin}

\gDefinition{Die Medizin ist die Lehre von Gesundheit und
Krankheit sowie die Kunde der Gesunderhaltung, Heilung und
Linderung.}

% \paragraph*{Einteilung} Notfälle kann man mehr oder weniger willkürlich
% einteilen
% 
% \begin{itemize}
% 	\item Störungen des Bewu{\ss}tseins --  neurologische Notfälle
% 	\item pulmonale Notfälle
% 	\item Kreislaufstörungen
% 	\item Gefä{\ss}verschlüsse
% 	\item Abdominelle Notfälle
% 	\item Urologische Erkrankungen
% \end{itemize}

\subsection{Begriffe}

Die Kenntnis von Fachbegriffen in der Medizin und
Sanitätshilfe ist nicht deshalb wichtig weil die Fachwörter
so gut klingen. Vielmehr dient die Fachsprache dazu daß
verschiedene Personen miteinander kommunizieren können und
dabei sichergestellt ist daß sie sich gegenseitig verstehen.

\paragraph{Symptome} Krankheitserscheinungen
\paragraph{Komplikationen} möglicherweise auftretende
Schwierigkeiten
\paragraph{Therapie} Ma{\ss}nahmen zur Heilung oder
Linderung von Krankheiten und Symptomen
\paragraph{Diagnose}
\begin{description}
    \item[Verdachtsdiagnose]
    \item[Notfalldiagnose]
    \item[Arbeitsdiagnose]
    \item['Status post' -- 'Zustand nach'] Abk. \gE{'St.p.'},
    \gE{'Z.n.'}. Bezeichnet eine ``alte Diagnose'', d.h. eine
    Erkrankung oder eine Behandlung die einmal durchgemacht
    wurde. Z.B.:
    \begin{itemize}
        \item \gTt{Z.n. Blinddarm-Operation}: Der Patient
        hatte einmal eine Blinddarm-Operation.
        \item \gTt{Z.n. Herzinfarkt 2/2006}: Der Patient
        hatte im Februar 2006 einen Herzinfarkt.
    \end{itemize}
\end{description}
\paragraph{Differentialdiagnosen (DD)} Diagnosen,
welche bei den jeweiligen Symptomen auch
möglich, aber nicht so wahrscheinlich sind wie jene
Diagnose für welche man sich ``entscheidet''.

Beispiel: Patient mit Kopfschmerzen nach ausgiebigen
Alkoholgenuß am Vortag.
\begin{description}
    \item[Verdachtsdiagnose:] \gE{``Kater''}
    \item[Differentialdiagnosen:]
    Migraine, Hirnblutung, Schlaganfall,
    Hypoglykämie, Intoxikation, \ldots
\end{description}
\paragraph{Diagnostik}
\paragraph{Befund}

\subsection{Medizinische Fachbereiche}

\subsection{Andere Medizinische Berufe --- Ein kleines
\emph{"`who is who"'}}

\subsection{Dokumentation -- Von der medizinischen Seite
betrachtet}

Ihre Dokumentation muß eine Geschichte erzählen. Eine
Geschichte die ein Dritter, unbeteiligter, verstehen
\gE{muß}. Dieser unbeteiligte Dritte muß allein anhand der
Schilderung in Ihrer Dokumentation fähig sein eine
Verdachtsdiagnose zu stellen. Weiters muss er
nachvollziehen können was mit dem Patient passiert ist.
Auch die Ereignisse davor oder die Einsatzsituation soll
für denjenigen verständlich sein (Unfallhergang,
Wohnsituation bei psychiatrischen Patienten, \ldots)

Schlecht:
\begin{quote}

{\sffamily Diagnose:} {\ttfamily RQW}

{\sffamily Befunde:} {\sffamily RR} {\ttfamily 130/70,}
{\sffamily HF} {\ttfamily 90} {\sffamily SpO2}
{\ttfamily 98\% }

{\sffamily Anmerkungen: --}
\end{quote}

Besser:

\begin{quote}
{\sffamily Diagnose:} {\ttfamily RQW Stirn}

{\sffamily Befunde:} {\sffamily RR} {\ttfamily 130/70,}
{\sffamily HF} {\ttfamily 90} {\sffamily SpO2}
{\ttfamily 98\% }

{\sffamily Anmerkungen:} {\ttfamily Pat. fuhr mit Fahrrad
und prallte mit Stirn gegen Ast., kein Sturz. Lt. Begleiter
keine B.losigkeit \underline{Traumacheck:} RQW Stirn, sonst unauffällig.
\underline{Neuro:} grob neurologisch unauffällig, Bew.
klar, Pat zeitl, örtl, zur Situation \& Person orientiert, Pupillen
mittelweit, Reakt. prompt und seitengleich, Kraft\&Gefühl
seitengl. in OE und UE. \underline{Anamn.:} \underline{Sympt:} Kopfschmerz seit
Aufprall. \underline{Krankh.:} keine bek. \underline{All.}
keine bek. \underline{Med.:} Aspirin heute früh. letzte
Tetanus-Impf. nicht bek. \underline{Ereignisse} s.o.
\underline{Letzte} Mahlzeit: heute Früh 2 Semmel

\underline{Maßnahmen:} Steriler Verband, nüchtern belassen,
Transport liegend mit erh OK.}
\end{quote}

Psychiatrisch:

\begin{quote}
{\sffamily Diagnose:} {\ttfamily Psychiatrische Erkankung,
V.a. Wahnvorstellungen, Unterbringung nach UbG}

{\sffamily Befunde:} {\sffamily RR} {\ttfamily 130/70,}
{\sffamily HF} {\ttfamily 90} {\sffamily SpO2}
{\ttfamily 98\% }

{\sffamily Anmerkungen:} {\ttfamily Pat. sitzend in Wohnung angetr., von Pol. mit Handschellen
gefesselt. Pat. warf lt. Polizei Gegstände aus d Fenster
auf Straße, verletzte Passanten ($\rightarrow$ FLO-1). Pat.
gibt an sich geg. Geheimdienst gewehrt zu haben,
jetzt agitiert und unkooperativ. 

\underline{Traumacheck:} keine äußeren Anz f Verletzungen,
sonst nicht erhhebbar (unkoop.). \underline{Neuro:} Pat.
ist wach, sonst nicht erhebbar

\underline{Anamn.:} \underline{Sympt:} keine Antwort
\underline{Krankh.:} lt. Spitalbrief Otto-Wagner-Sp. v.
1.4.08 paranoide Schizophrenie \underline{All.} k.A.
\underline{Med.:} nicht eruierbar, auf Tisch mehrere
angebr. Packungen Rahypnol \underline{Ereignisse} s.o.
\underline{Letzter Spitalsaufenthalt} k.A.

\underline{Wohnsituation:} heruntergekommen, vermüllt, Pat.
wohnt mit großem Hund, ganze Wohnung verkotet, bestialisch
stinkend. Auf Tisch mehrere Schachteln "`Rohypnol"'.

\underline{Maßnahmen:} Einweisung nach UbG, Transport in
Handschellen in Pol-begl. (Viktor 6, DNr. 4711). Protokoll
Amtsarzt ad. Spital.}
\end{quote}


\section{Herz-Kreislaufstörungen}

\dictum[Goethe]{Blut ist ein ganz besond'rer Saft.}

\label{bkm:RefHeading42342802}

\paragraph{Überblick über die wichtgsten
Herz-Kreislauf-Notfälle}

\begin{itemize}
    \item Herz-/Kreislaufstillstand
    \item Herzinsuffizienz
    \begin{itemize}
        \item Linksherz-Versagen
        \item Rechtsherz-Versagen
    \end{itemize}
    \item (Akutes) Koronarsyndrom
    \begin{itemize}
        \item Angina pectoris
        \item Myokardinfarkt (MCI )
    \end{itemize}
    \item Herzrhythmusst\"orungen
    \item Störungen des Bludrucks
    \begin{itemize}
        \item Hyp\gEs{o}tonie
        \item Hyp\gEs{er}tonie
        \begin{itemize}
            \item Hypertensive Krise
            \item Hypertensiver Notfall
        \end{itemize}
    \end{itemize}
\end{itemize}


\subsection{Herz-Kreislauf-Stillstand}

\marginline{Reanimation}Der Herz-Kreislauf-Stillstand wird
im Kapitel Reanimation besprochen.

\subsection{Herzinsuffizienz}

\gDefinition{Unzureichende Auswurfleistung des Herzens.}

Leistungseinschränkung des Herzmuskels f\"uhrt zu einer verminderten
Auswurfleistung, viele Ursachen sind m\"oglich. Manche Ursachen treten
\gE{akut} auf, manche sind eher \gE{chronischer} Natur.

Je nachdem welcher Herzteil betroffen ist unterteilt man in:

\begin{description}
	\item[\gT{Rechtsherzinsuffizienz}]\index{Rechtsherzinsuffizienz} Unzureichende Leistung des rechten Herzens, 
	\item[\gT{Linksherzinsuffizienz} ]\index{Linksherzinsuffizienz} Unzureichende Leistung des lingen Herzens, und
	\item[\gT{Globalinsuffizienz},]\index{Globalinsuffizienz} Unzureichende Leistung beider Herzteile.  
\end{description}

Abhängig davon wie gut der Körper mit der Einschränkung zurecht kommt unterteilt man: 

\begin{description}
	\item[Kompensierte Herzinsuffizienz] Der K\"orper hat noch genug Reserven, kann gegensteuern und kommt mit der eingeschränkten Leistung aus.
	\item[Dekompensierte Herzinsuffizienz] Der K\"orper kommt aufgrund des Fortschreitens der Erkrankung oder aufgrund eines erh\"ohten Bedarfs nicht mehr mit der Leistung aus. Es treten deutliche Symptome auf. 
\end{description}

\paragraph*{Beurteilung des Schweregrades}

Es ist wichtig den Schweregrad einer Herzinsuffizienz zu beurteilen. Sie
muss nicht dramatisch sein, f\"ur viele Patienten ist sie ein
chronischer Zustand mit dem sie gut leben k\"onnen.
Problematisch wird es jedoch wenn der Patient pl\"otzlich eine Symptomatik entwickelt und
die Kompensationsmechanismen (k\"opereigene, Medikamente, {\dots})
versagen. Man spricht dann von der \gE{dekompensierten} Herzinsuffizienz.
Man muss die Symptome generell \gE{anhand ihres
Verlaufes beurteilen}: Bekannte, seit längerer Zeit bestehende
Symptome sind weniger wichtig als pl\"otzlich auftretnde Symptome,
diese sind meistens \gEs{Alarmzeichen}. 

Weitere klassiche Alamrzeichen sind Symptome wie
\gE{Atemnot}, \gE{pl\"otzlich auftretende
\"Odeme} oder \gE{akute Schmerzen}. Ein spezielles,
hochgefährliches, Alarmzeichen ist das
\gEs{brodelnde Atemgeräusch} beim Lungen\"odem. 

Das kardiale Lungen\"odem wird im Kapitel \ref{topic:lungenoedem} (Lungen\"odem)
behandelt. 


\subsubsection{Linksherzinsuffizienz}

\index{Linksherzinsuffizienz}
Die Linksherzinsuffizienz entsteht aufgrund mangelnder Auswurfleistung
des linken Herzens. Durch den R\"uckstau von Blut in den
Lungenkreislauf kann es zu einem kardialen Lungen\"odem kommen. 

Patienten mit einem \gE{kardialem Lungen\"odem} sind \gE{lebensbedrohlich} krank.
Das Herz vebraucht oft schon die allerletzten Reserven, schon die
Aufregung durch den Transport im Stiegenhaus zum Auto kann das Herz
endg\"ultig in den Herz-Kreislaufstillstand bringen. 

Es kommt immer wieder vor das solche Patienten im Stiegenhaus oder im
Auto reanimationspflichtig werden. 

\gCave{Kein Transport von Patienten mit Lungen\"odem ohne Notarzt. Auch wenn
das Spital {\quotedblbase}um{\textquotesingle}s Eck{\textquotedblleft}
ist.}

\paragraph{Symptome einer Linksherzinsuffizienz}

\begin{itemize}
    \item {\mdseries
    \gEs{Dyspnoe} !}
    
    \begin{itemize}
        \item Orthopnoe: Patient bekommt nur mehr im Sitzen Luft, kann nicht flach
        schlafen, etc.
    \end{itemize}
    \item \gEs{Tachypnoe}
    \item \gEs{Tachykardie}
    \item \gEs{Hustenreiz} !
    \item Unruhe
    \item \gEs{Hypotonie} !
    \item \gEs{Lungenstauung/Lungen\"odem} (Blut staut sich
    zur\"uck in den Lungenkreislauf)
\end{itemize}


\begin{gMassnahmenEnv}
    \minisec{Spezielle Ma{\ss}nahmen}
    
    \gTxBeiVit \gTxMassVitBes
    
    \begin{itemize}
        \item Lagerung: Oberk\"orper hoch
        \item $O_2$ 10~l, bei Bedarf bis zu 15~l
        \item Ursachenforschung
    \end{itemize}
    
    sonst:
    
    \begin{itemize}
        \item Lagerung: Oberk\"orper hoch
        \item $O_2$ je nach Symptomen 6-8~l, bei Bedarf bis zu 15~l
        \item Vitalwerte erheben \& Monitoring
    \end{itemize}
\end{gMassnahmenEnv}


\subsubsection{Rechtsherzinsuffizienz}
Die Rechtsherzinsuffizienz entsteht aufgrund mangelnder Auswurfleistung
des rechten Herzens. 

\subparagraph{Symtome}

% \begin{wrapfigure}{r}
% %	\begin{figure}
% 		%\caption{Aszites und Bein"odeme}	
% 		\includegraphics[width=5cm]{./images/aszites-beinoedeme.png}
% %	\end{figure}
% \end{wrapfigure}

\begin{itemize}
	\item \gEs{Bein-}/\gEs{Haut\"odeme}
	\item Nächtliches Wasserlassen
	\item \gEs{gestaute Halsvenen} 
	\item Lebervergr\"o{\ss}erung, evt. mit Aszites
	({\quotedblbase}Wasserbauch{\textquotedblleft})
\end{itemize}

\begin{gMassnahmenEnv}
    
    \paragraph{Spezielle Ma{\ss}nahmen} ~
    
    \gTxBeiVit \gTxMassVitBes
    \begin{itemize}
        \item Lagerung: Oberk\"orper hoch
        \item $O_2$ 10~l, bei Bedarf bis zu 15~l
    \end{itemize}
    
    sonst:
    \begin{itemize}
        \item Lagerung: Oberk\"orper hoch
        \item $O_2$ je nach Symptomen 6-8~l, bei Bedarf bis zu 15~l
        \item Vitalwerte erheben \& Monitoring
    \end{itemize}
    
    weiters:
    \begin{itemize}
        \item Ursachenforschung
    \end{itemize}
    
\end{gMassnahmenEnv}

\subsection{Koronare Herzkrankheit und Akutes Koronarsyndrom}



\begin{description}
	\item[Die \index{Koronare Herzkrankheit}\gT{Koronare Herzkrankheit}
	(\index{KHK}\gT{KHK})] ist eine Erkrankung der
	\gEs{Herzkranzgefä{\ss}e} (Koronargefä{\ss}e)
	bei der es zu einer Verengung der Gefä{\ss}e kommt.

	\item[Das \index{Akutes Koronarsyndrom}\gT{Akute Koronarsyndrom}] ist ein
	Sammelbegriff f\"ur \gE{Herzinfarkt}, \gE{stabile} und \gE{instabile Angina pectoris}.
	Bei diesen Krankheitsbildern kommt es zu einer akuten
	\gEs{Unterversorgung des Herzmuskels} mit Blut und
	Sauerstoff. Die Ursache sind in erster Linie verengte oder verstopfte
	Herzkranzgefä{\ss}e . 

\end{description}

\subsubsection[Angina pectoris]{Angina pectoris}
\index{Angina pectoris}{\mdseries
Pl\"otzliches Engegef\"uhl in der Brust meist mit Thoraxschmerz infolge
einer Unter\-ver\-sorgung des Herzmuskels mit sauerstoffreichen Blut.
Der Anfall vergeht wieder von selbst ohne bleibende Schäden.}

\begin{itemize}
	\item Minderdurchblutung des Herzmuskels (Myokard) aufgrund einer oder
	mehrerer Eineingung(en) von Herzkranzgefä{\ss}ren (Koronararterien)
	oder
	\item Erh\"ohter Sauerstoffbedarf des Herzens bei Anstregung
\end{itemize}
{\mdseries
${\rightarrow}$ Mi{\ss}verhältnis Sauerstoffangebot - 
Sauerstoffbedarf des Herzens ${\rightarrow}$ Schmerz!}

Der akute Angina-Pectoris-Anfall geh\"ort zu der Familie des Akuten
Koronarsyndroms. Dazu geh\"ort auch der Herzinfarkt, dem der gleiche
Mechanismus wie der der Angina Pectoris zu Grunde liegt, nur mit
schwer\-wiegenderen Auswirkungen. 

Bei den meisten Patienten ist die Angina Pectoris schon bekannt, bzw. es
besteht eine {\quotedblbase}\index{Koronare
Herz-Krankheit}\gT{Koronare Herz-Krankheit}
(\index{KHK}\gT{KHK}){\textquotedblleft}. Die KHK ist
eine Erkrankung der Herzkranzgefä{\ss}e, diese werden durch
Ablagerung von Fetten und Kalk geschädigt und eingeengt, dadurch kann
weniger Blut durchflie{\ss}en. So kommt es zu einer Mangelversorgung
des Herzmuskels und zu den Symptomen der Angina pectoris (oder zu einem
Herzinfarkt). 


\begin{center}
 [Warning: Image ignored] {AASdev20090228-img72.png}
\captionof{figure}{\gAuthNotes{Bild fehlt!} Patient mit
einem akuten Koronarsyndrom, vermutlich Angina Pectoris. Bei kaltem Wetter, beim
Stiegen steigen mit schwerem Gepäck: Ein pl\"otzliches Engegef\"uhl
im Brustkorb und ein stechender Schmerz im linken Thorax.}

\end{center}

\paragraph{Symptome}

\begin{itemize}
	\item Brustschmerz mit Ausstrahlung
	\item Todesangst
	\item Schwächegef\"uhl
	\item Dyspnoe
	\item primär \gEs{nicht sicher von MCI unterscheidbar} 
	\item Symptome treten eher bei Belastung auf und vergehen oft bei Entspannung.
\end{itemize}

\begin{center}
 [Warning: Image ignored] {AASdev20090228-img73.png}
\captionof{figure}{\gAuthNotes{Bild fehlt!}
Schmerzausdehnung beim ischämischen Herzschmerz.}


\end{center}

\begin{gMassnahmenEnv}
\paragraph{Spezielle Ma{\ss}nahmen}

\begin{itemize}
	\item \gTxMassVitBes
	\begin{itemize}
		\item Lagerung: Oberk\"orper hoch
		\item $O_2$ 10~l, bei Bedarf bis zu 15~l
	\end{itemize}
	\item \gEs{striktes Bewegungsverbot}!
	\item Pat. darf eigenes Nitro nur  bei RR {\textgreater}100 syst. nehmen!
\end{itemize}
\gAuthNotes{
??? eigenes Nitro nehmen?

Vorschlag: Nur wenn sein Nitro leer war und er es deswegen nicht
genommen hat bei RRsyst{\textgreater}110

Wenn er es schon genommen hat und es nicht geholfen hat
dann nicht}
\end{gMassnahmenEnv}

\subsubsection{Herzinfarkt - Myokardinfarkt (MCI)}

\gDefinition{Verschluß eines oder mehrerer Herzkranzgefäße mit
anschlie{\ss}endem Absterben des betroffenen Herzmuskelgewebes.}

\gMarried{\paragraph{Symptome}

Die Symptome entsprechen im Wesentlichen denen der Angina pectoris, doch
vergehen diese nicht mehr und es bleibt auch bei optimaler Behandlung
ein Schaden am Herzmuskel. Bei Diabetikern kann es aufgrund der
Nervenschädigungen auch zu einem schmerzlosen Infarkt kommen, man
spricht dabei auch vom \gT{{\quotedblbase}stummen
Infarkt{\textquotedblleft}}\index{Stummer Infarkt}.

Sonst steht normalerweise der Brustschmerz zusammen mit
Atemnot und Todesangst im Vordergrund. Der
\gEs{Brustschmerz} kann in Richtung Magen, rechter oder linker Hand, Kiefer
oder auch Rücken \gEs{ausstrahlen}. Zusätzlich zu dem
Schmerz verspürt der Patient auch oft ein \gEs{Engegefühl} im
Brustkorb, \emph{"`als würde jemand darauf stehen"'}. Die
Atemnot kann massiv ausfallen. 

Weiters können aufgrund des eintretenden \gEs{kardiogenen
Schocks} die entsprechenden Schockzeichen beobachtet
werden: Blässe, Kaltschweißigkeit, Hypotonie, etc.

\subparagraph*{Komplikationen -- Womit muss ich rechnen?}

\begin{itemize}
    \item Herzrhythmusst\"orungen (alle Arten), bis hin
    zum Atem-/Kreislaufstillstand und Reanimation
    \item Kardiogener Schock
\end{itemize}}{
\paragraph{Schlagworte}

\begin{itemize}
\item Verschluß von Herzkranzgefäßen
\item Absterben von betroffenen Herzmuskelgewebes
\item Symptome wie bei Angina Pectoris, vergehen aber nicht
mehr
\begin{itemize}
	\item Brustschmerz mit Ausstrahlung 
	\item Druckgef\"uhl / Beklemmung {\quotedblbase}als w\"urde jemand am
	Brustkorb stehen{\textquotedblleft}
	\item Todesangst
	\item Dyspnoe
	\item Kardiogener Schock
	\item Keine ausreichende Besserung auf Nitro
\end{itemize}
\item Bei Diabetikern schmerzloser Infarkt möglich
(Nervenschädigung))
\item \gEs{Lebensgefahr!}
\item Jede Rhythmusstörung als Komplikation möglich.
\item Behandlung durch
\begin{itemize}
    \item Lyse durch Notarzt
    \item Herzkatheter (PTCA, PCI) im Spital
\end{itemize}
\end{itemize}}

\begin{center}
 
 \gTempPicMissing{Infarkt [Warning: Image ignored]
 {AASdev20090228-img74.png}}

\end{center}
\begin{gMassnahmenEnv}
    \paragraph{Spzielle Ma{\ss}nahmen}
    
    \begin{itemize}
        \item \gTxMassVitBes 
        \begin{itemize}
            \item Lagerung: Oberk\"orper hoch
            \item $O_2$ 10~l, bei Bedarf bis zu 15~l
        \end{itemize}
        \item Beengende Kleidung \"offnen
        \item Striktes Bewegungsverbot!
    \end{itemize}
\end{gMassnahmenEnv}

\subparagraph{Weitere Therapie}

Herzinfarktpatienten werden {\quotedblbase}besonders{\textquotedblleft}
behandelt und normalerweise nicht auf eine beliebige Interne Station
gebracht. Meistens stehen dem Notarzt zur Auswahl:

\begin{itemize}
	\item Lyse-Therapie mittels eines speziellen Medikaments, oder
	\item Einweisung in ein \index{Herzkatheterlabor}\gEs{Herzkatheterlabor}
	({\quotedblbase}PCI{\textquotedblleft},
	{\quotedblbase}\index{PTCA}\gEs{PTCA}{\textquotedblleft}
	\gZ{{}- Perkutane Transluminale Coronar
	Angioplastie})

\end{itemize}

In Wien gibt es einen \gE{Dienstplan} f\"ur die Katheterlabore, d.h. 
jedes Katheterlabor hat an bestimmten Tagen der Woche rund um die 
Uhr Bereitschaftsdienst. Eine Voranmeldung und vorangehende Absprache 
zwischen Notarzt und dienstahbenden Oberarzt der Abteilung ist erforderlich. 
Die PTCA ist in Wien u.a. verf\"ugbar im AKH, Donauspital, und KH Hietzing.

\gCave{{\bfseries Kein Transport ohne Notarzt!}
	
	Auch nicht wenn
	das Spital {\quotedblbase}um{\textquotesingle}s Eck{\textquotedblleft}
	ist.

	Der Patient braucht vermutlich {\bfseries kein beliebiges
	Spital} sondern ein {\bfseries dienst\-bereites
	Herzkatheterlabor}. Der Notarzt muss entscheiden ob das notwendig und sinnvoll ist. 

	Auch wenn der Patient ein Einweisungsformular (Spitalszettel) von einem	
	Arzt hat (\"Arztefunkdient, Hausarzt) oder ein Hausarzt meint
	{\quotedblbase}es sei eh nicht dringend{\textquotedblleft} unbedingt
	einen Notarzt beiziehen.
}

\subsection{Herzrhythmusst\"orungen }

\begin{center}
\begin{table}[h]
\captionabove{Wichtige und häufige Herzrhythmusstörungen}
\label{tab:herzrhythmusstoerungen}
\begin{tabulary}{\textwidth}{LLL}
\toprule
Rhythmus & Bedeutung & Bemerkung
\\\midrule

\gEs{Bradykardie} 

{\small (HF\,{\textless}\,60/min)}\footnote{beim
Erwachsenen} & Situations\-abhängig & \emph{"`zu langsam"'}
\\
\gEs{Tachykardie} 

{\small (HF\,{\textgreater}\,100/min)}\footnote{beim
Erwachsenen} & Situations\-abhängig & \emph{"`zu schnell"'}
\\\midrule \gEs{Kammerflimmern} & Reanimation  & ?
\\
\gEs{Pulslose Ventrikuläre Tachy\-kardie} & Reanimation & ?
\\
\gEs{Asystolie} 

{\small (HF\,=\,0)} & Reanimation & ?
\\
\gEs{Vorhofflimmern} & regellose Erregung im Vorhof

oft symptomlos & 
\gE{Chronische Erkrankung.} Sehr viele Leute haben
Vorhofflimmern und leben ganz gut damit. 

Risikofaktor für:

\gE{Tachykardie-Anfälle}

\gE{Thrombenbildung}, dadurch erhöhtes Lungenembolie-
und Schlaganfallrisiko. Deshalb Vorbeugung mit
gerinnungs\-hemmenden ({\quotedblbase}blut\-ver\-d\"unnenden{\textquotedblleft}) Medi\-ka\-menten:
z.B. \gE{"`Marcoumar"'} ${\rightarrow}$ Patient ist blutungsgefährdet
\\\bottomrule
\end{tabulary}
\end{table}
\end{center}



\subsubsection{Hier zählt die Geschwindigkeit: Tachy- und
Bradykardie} 

\begin{description}
    \item[Tachykardie] \ldots zu schneller
    \marginpar{Tachometer im Auto:
    Geschwindigkeitsanzeiger, meistens zu schnell ;)}
    Herzschlag
    \item[Bradykardie] \ldots zu langsamer Herzschlag
\end{description}

Denke bei der Beurteilung auch immer an die Umstände und den
Erregungszustand des Patienten!

\subparagraph{Bradykardie}
HF {\textless} 60/min (beim Erwachsenen)

Leistungssportler k\"onnen augrund ihres trainierten K\"orpers auch sehr
niedrige Ruhepulsfrequenzen haben. 

\subparagraph{Tachykardie }
HF {\textgreater}100/min (beim Erwachsenen)

Tachykarde Attacken sind ein häufiges Notfallbild. Da oft chronische
Erkrankungen die Ursache sind  haben die Patienten diese Attacken immer
wieder und kennen die Symptome bereits (und oft auch die Behandlung). 

\subparagraph*{Symptome}

\begin{itemize}
    \item {\mdseries
    Herzklopfen}
    
    \item {\mdseries
    evt. Leichter Brustschmerz}
    
\end{itemize}

\paragraph{Gefahr}

Da das \gEs{Herz pl\"otzlich viel mehr Arbeit leisten
muss} kann es sein dass es an seine \gEs{Grenzen}
st\"osst und \gEs{versagt} bzw. die Blutversorgung
des Herzens f\"ur diese Arbeit nicht mehr reicht und ein Herzinfarkt
entsteht.

\begin{gMassnahmenEnv}
    \subparagraph*{Spezielle Ma{\ss}nahmen bei symptomatischer Tachykardie}
    
    \begin{itemize}
        \item Bei vitaler Bedrohung: Standardma{\ss}nahmen bei vital bedrohten
        Patienten. Besonderheiten: 
        \begin{itemize}
            \item Lagerung: Oberk\"orper hoch
            \item O\gSub{2} 10~l, bei Bedarf bis zu 15~l
        \end{itemize}
        \item Beengende Kleidung \"offnen
        \item Striktes Bewegungsverbot!
        \item Psychische Betreuung! Oase der Ruhe schaffen. 
    \end{itemize}

\end{gMassnahmenEnv}



\subsubsection{Besondere Rhythmen}

\paragraph[Kammerflimmern]{Kammerflimmern}

\gEs{Reanimationspflichtig}. NICHT verwchseln mit Vorhofflimmern! Die
elektrische Erregung im Herzen ist ungerichtet, als w\"urden alle Autos 
eines Parkplatzes gleichzeitig in alle Richtungen losfahren 
${\rightarrow}$ \gEs{Kein Auswurfleistung} 
    

\paragraph{Pulslose Ventrikuläre Tachkardie}
\gEs{Reanimationspflichtig!} Vorstufe zum Kammerflimmern, Herz kontrahiert so
schnell dass es sich nicht mehr f\"ullen kann ${\rightarrow}$ \gEs{keine Auswurfleistung} 
    
\paragraph{Asystolie}
\gEs{Reanimationspflichtig}. Im Herz entsteht keine elektrische
Erregung$\rightarrow$  \gEs{keine Auswurfleistung} 
    
\paragraph[Vorhofflimmern]{Vorhofflimmern}

Regellose Erregung im Vorhof, die oft symptomlos ist. 
Sehr viele Leute haben Vorhofflimmern und leben ganz gut damit.

Komplikationen:
\begin{itemize}
    \item Tachykardie-Anfälle
    \item Thrombenbildung
    \begin{itemize}
        \item deshalb Therapie mit \gEs{Gerinnungshemmenden}
        ({\quotedblbase}blutverd\"unnenden{\textquotedblleft}) Medikamenten:
        Marcoumar ${\rightarrow}$ Patient ist Blutungsgefährdet
    \end{itemize}
\end{itemize}

\subsection{Hypotonie}

\gDefinition{Unangemessen zu niedriger Blutdruck.}
 
\gEs{Richtwert:} Beim Erwachsenen RR {\textless} 100 mmHg
syst. Jeder Mensch hat allerdings im gewisen Rahmen einen
"`eigenen Normalwert"', bei Kindern gelten grundsätzlich
andere Werte.


\paragraph{Ursachen}

\begin{itemize}
    \item Cardial
    \item Volumenmangel, Exsikkose
    \item Regulationsst\"orung (vasovagal)
    \item Sportler, junge Frauen
    \item Schock
\end{itemize}

\paragraph{Symptome}

\begin{itemize}
    \item M\"udigkeit bis Kollaps
    \item Kollaps
    \item RR-Abfall mit Bewu{\ss}tseinseintr\"ubung
    \item {\quotedblbase}schwarzwerden vor Augen{\textquotedblleft}, Ohnmacht
    \item Sturz = {\quotedblbase}Therapie{\textquotedblleft}
\end{itemize}


\subsubsection{Kollaps, Synkope}

\gDefinition{passagere Kreislaufregulationsst\"orung  mit RR-Abfall,
    Blutverteilungsst\"orung  ohne Volumsverlust}

\paragraph{Symptome}

\begin{itemize}
    \item RR ${\downarrow}$
    \item Schwindel (Vertigo )
    \item evt. kurze Bewu{\ss}tlosigkeit, Ohnmacht
    \item Patient sackt zusammen
    \item Bradykardie oder Tachykardie
    \item Schwarzwerden vor Augen
\end{itemize}

Patient wacht auf dem Boden liegend  auf, wei{\ss} nicht, was geschehen
ist...

\paragraph{Ursachen} Die Ursachen sind präklinisch oft
nicht sicher ermittelbar, mitunter bleibt die wahre Ursache
für immer ein Rätsel. Daher ist eine Hospitalisierung zur
bestmöglichen Abklärung angeraten. 

Es gibt jedoch auch typische "`Kollaps-Situationen"': Ein
schwüler Tag in einer überfüllten U-Bahn-Garnitur, nicht
gefrühstückt, auf dem Weg zu einem Vorstellungstermin und
seit gestern in der Menstruation. Oder im stickigen
Veranstaltungszelt bei 40\textdegree , wobei der der ganze
Flüssigkeitskonsum aus 5 (entwässerden) Dosen
RedBull\texttrademark\footnote{Achtung! Schleichwerbung!}
bestand. So mancher war auch ob eines Begräbnisses so
ergriffen daß er sich fast in die Grube neben den Sarg
dazugelegt hätte \ldots 

ausgehend von:

\begin{itemize}
    \item Herz\footnote{z.B. Herzrhythmusstörungen}
    \item Gefäßen\footnote{z.B. Hitzebedingte Erweiterung
    der GEfäße mit "`Versacken"' des Blutes in den Beinen}
    \item Blut\footnote{z.B. Blutarmut (\gT{Anämie}),
    Regelblutung}
    \item Hirn
    \item Hirngefäßen\footnote{z.B. TIA, s.\,\ref{topic:tia}
    S.\,\pageref{topic:tia}}
\end{itemize}

\gAuthNotes{ERKL\"AREN!}

\begin{gMassnahmenEnv}
    \paragraph{Spezielle Maßnahmen bei Kollaps/Synkope}
    
    \begin{itemize}
        \item Erheben ob eine akute vitale Bedrohung besteht
        \item Ausschluß von Notfallsdiagnosen soweit
        möglich, die gesammten zur Verfügung stehenden
        Möglichkeiten sind auszuschöpfen.
        \item Im Besonderen:
        \begin{itemize}
            \item Traumacheck (Sturz?)
            \item Neurocheck inkl. Blutzuckermessung
            \item Temperatur
            \item Ausführliche Anamnese

        \end{itemize}
        
        \item Lagerung: Nach Ausschluß von Herzbeschwerden,
        Atemnot und relevanten Verletzungen: Beine hoch. 
        \item Hospitalisierung zur Abklärung
        anstreben\newline
        Abteilung: Innere Medizin
        
    \end{itemize}
    
\end{gMassnahmenEnv}

\subsection{Hypertonie}

\index{Hypertonie}\gT{Hypertonie}
(\index{Bluthochdruck}\gT{Bluthochdruck}): Oft
symptomlose \textit{chronische} Erkrankung welche auf Dauer schwere
Schäden am K\"orper anrichten kann und im Zuge von pl\"otzlichen
Blutdruckkrisen auch akut Probleme machen kann. 

\begin{itemize}
    \item Chronische Hypertonie: Grenze derzeit bei 140/90
    (au{\ss}er bei Kindern)
\end{itemize}

\begin{gCaveEnv}
    Bei Schwangeren sind die Grenzwerte peinlich genau zu beachten! Denke an
    die \gT{EPH-Gestose}\footnote{Die EPH-Gestose wird in
    Kap. \ref{topic:eph-gestose}
    auf Seite \pageref{topic:eph-gestose} erklärt.}!
\end{gCaveEnv}

\paragraph{Ursachen}

\begin{itemize}
    \item Widerstandserh\"ohung (durch Engstellung der Gefä{\ss}e)
    \item vermehrte Herzarbeit
    \item Sympathikusaktivierung, Adrenalin+Cortison (Stre{\ss}hormone) durch
    Dauerstre{\ss},  Bewegungsmangel, genetische Veranlagung  
\end{itemize}

\subparagraph{Langzeitfolgen}

Die Hypertonie hat in der Notfallmedizin auch eine wichtige
indirekte Bedeutung: Sie sit Ursache für viele andere
Erkrankung, die ihrerseits akut lebensbedrohlich werden
können. 

Zwei Wirkungen der Hypertonei sind dabei besonders wichtig:
Die Gefäßschödigung und die Bealastung des Herzens durch
den hohen Gegendruck. 

Bei der Gefäßschödigung kommt es zur ABlagerung durch Kalk
und sonstige Plaques udn über längere Zeit zur Verengung
oder Verstopfung der Gefäße. Je nach Organ kann das sehr
massvie Auswirkungen haben. Im Herz führt es zur Angina
Pectoris oder zum Herzinfakr, im Hirn zu Schlaganfällen, in
den Nieren zu eine Niereninsuffizienz, usw. Selbst große
Gefäße wie die Hauptschlagader (Aorta) können derartig
geschädigt werden dass sie reißen (Aortenaneurysma) --- eine
für den Patienten sehr "`unangenehme"' Sache. 

Die Schödigung des Herzmuskels ist die Folge von dem
erhöhten Durck gegen den das Herz pumpen muß. Der Muskel
vergrößert sich um mit der Belastung fertig zu werden, kann
aber damit nicht mehr ausreichend mit sauerstoffreichen.
Blut versorgt werden. 

Wie haben also gesehen daß die Hypertonie ein wichtger
Risikofaktor für viele schwere Notfälle ist. Daher ist es
auch verständlich daß die Behandlung der Hypertonie einen
großen Stellenwert in der medizinischen Versorgung der
Bevölkerung hat. 

\begin{itemize}
    \item Gefä{\ss}schädigung
    \begin{itemize}
        \item Herzkranzgefä{\ss}e ${\rightarrow}$ Infarkt
        \item Hirngefä{\ss}e ${\rightarrow}$ Schlaganfall
        \item Nierengefä{\ss}e ${\rightarrow}$ Niereninsuffizienz ${\rightarrow}$
        Dialyse ${\rightarrow}$ typischer KTW-Patient
        \item Hauptschlagader ${\rightarrow}$ Aneurysma        
    \end{itemize}
    \item Herzmuskel (Herz muss gegen gr\"o{\ss}eren Widerstand pumpen)
    \begin{itemize}
        \item dilatative Kardiomyopathie
    \end{itemize}
\end{itemize}

\subsubsection{Hypertensive Krise und Hypertensiver Notfall}

{\bfseries\color[rgb]{0.5019608,0.0,0.0} SEHR FRAGLICH !!!}

\begin{center}
\begin{tabularx}{\textwidth}{XX}
\toprule
Hypertensive Krise
 &
Hypertensiver Notfall
\\\midrule
{\mdseries RR {\textgreater} 230/130 mmHg}

{\mdseries Patient f\"uhlt sich gut}

 &
{\mdseries RR {\textgreater} 230/130 mmHg }

{\mdseries + Symptome}

\\\hline
{\mdseries Zufallsbefund?}

{\mdseries Oder steht der Blutdruck doch in Zusammenhang mit der
Berufung?}

{\mdseries \textit{Patient hat \"uberhaupt keine Beschwerden?} Warum hat
er Sie dann \"uberhaupt gerufen?}

{\mdseries Wenn es wirklich nur ein Zufallsbefund ist, zum Beispiel wenn
Sie {\quotedblbase}aus Spa{\ss}{\textquotedblleft} jemand Blutdruck
gemessen haben ({\quotedblbase}Gesundheitstage{\textquotedblleft},
{\quotedblbase}Vorsorgeuntersuchung{\textquotedblleft}):}

\paragraph{Ma{\ss}nahmen}

\begin{itemize}
    \item Lagerung: Oberk\"orper hoch
    
    \item evt. $O_2$
    
    \item Beengende Kleidung öffnen
    
    \item Monitoring
    
    \item  Hospitalisierung oder Notarzt zur Belassung
    
\end{itemize}
&
\subparagraph*{ Symptome:}

\begin{itemize}
    \item Kopfschmerz
    \item Sehst\"orungen
    \item Brustschmerz
    \item evt. Nasenbluten
    \item \"Ubelkeit, Erbrechen, Schwindel
\end{itemize}
\paragraph*{Gefahr}

\gEs{Herzversagen (Herz muss
pl\"otzlich gegen gro{\ss}en Druck pumpen)}

\paragraph{Ma{\ss}nahmen}

\begin{itemize}
\item Lagerung: Oberk\"orperhoch
\item O2 
\item Beengende Kleidung \"offnen
\item Bewegungsverbot
\item Notarzt
\item Monitoring
\item Reanimationsbereitschaft
\end{itemize}
\\\bottomrule
\end{tabularx}
\end{center}

\subsection{Gefä{\ss}verschlüsse in den Extremitäten}\marginpar{Gefäßverschlüsse}

%%%
\marginline{
\begin{itemize}
    \item \textbf{\textcolor[rgb]{0.7882353,0.0,0.08627451}{Arteriell}}
    \item \textbf{\textcolor{blue}{Ven\"os}}:
    \index{tVT}\gT{tVT} (tiefe Venen Thrombose)
\end{itemize}
}%%%marginline
%%%
Grundsätzlich unterscheidet man je nach betroffenen Gefäß
zwischen \gEs{arteriellen} und \gEs{venösen}
Gefäßverschlüssen. Beim \gE{arteriellen Verschluß} besteht
das Hauptproblem in der Unterversorgung der Extremität mit
Blut. Beim \gE{venösen Verschluß} kommt es zu einer Ab- und
Rückflußbehinderung. Zusätzlich besteht die Gefahr daß sich
ein Gerinnsel losreißt und in einem anderen Teil des
Körpers Schaden anrichtet\footnote{Zum Beispiel in der
Lunge im Rahmen einer Lungenembolie}.

\subsubsection{Arterieller
Gefä{\ss}verschlu{\ss}}\marginlabel{Arterieller Gefä{\ss}verschlu{\ss}}

Beim \gE{arteriellen Gefäßverschluß} verschließt etwas das
zuführende Gefäß, es kommt zu einer schmerzhaften 
Unterversorgung der Extremität mit Blut. Der komplette
Verschluß passiert plötzlich. 

\paragraph{Warum kommt es zu einem Gefäßverschluß?}
Meistens liegt eine 
\gMargin{Chronische Gefäßschädigung}chronische Schädigung der Gefäße infolge anderer Erkankungen oder aufgrund eines ungesunden
Lebensstils vor. Ursachen für solch eine Gefäßschädigung
sind z.B. die Zuckerkrankheit (\gT{Diabetes mellitus},
s.\ref{topic:diabetes}), Hypertonie, erhöhte Blutfette oder
das Rauchen. Die Gefäße verkalken und verändern sich auch sonst
nicht zu ihrem Vorteil.

Wenn die Gefäße nur verengt sind spricht man von
der \gMargin{\gT{periphere arterielle
Ver\-schluß\-krank\-heit}, Abkz. \gT{paVK}}\gT{peripheren
arteriellen Verschluß\-krankheit}, abgekürzt \gT{paVK}. Sie ist eine chronische
Erkankung die man oft in Arztbriefen oder auf
Transportscheinen lesen kann. (Achtung: Obwohl es
"`\emph{Verschluß}krankheit"' heißt sind die Gefäße nicht komplett verschlossen sondern nur verengt!)

\gMarried{\paragraph{Symptome} 
Die Symptome ergeben sich aus der Unterversorgung der
Extremität mit sauerstoffreichem Blut: sie ist blaß,
kalt, tut höllisch weh und die Rekap-Zeit ist nicht
mehr sinnvoll messbar. Zusätzlich wird die Extremität
gefühllos und gelähmt. Der Patient ein Kribbeln
empfinden (\gE{"`Ameisenlaufen"'}. Tabelle \ref{tab:5p}
listet die klassischen Symptome auf \gZ{(die \emph{"`5
englischen P's"'})}.
}{
{
\begin{center}
\begin{tabulary}{4cm}{CLL}\toprule
\multicolumn{3}{c}{\gEs{Arterieller Gefäßverschluss}}\\
\midrule\addlinespace
\centering 1. & \gZ{Pain}        & \gEs{schmerzend}
\\
\centering 2. & \gZ{Pale}        & \gEs{bla{\ss}}
\\
\centering 3. & \gZ{Pulsless}    & \gEs{pulslos}
\\
\centering 4. & \gZ{Paresthesia} & \gEs{gef\"uhllos}
\\
\centering 5. & \gZ{Paralysis}   & \gEs{gelähmt}
\\\addlinespace\bottomrule
\end{tabulary}
\end{center}}
\captionof{table}{Die 5 Zeichen des arteriellen Gefäßverschluss.}\label{tab:5p}
}



\begin{gMassnahmenEnv}
    \paragraph{Spezielle Ma{\ss}nahmen}
    \begin{itemize}
        \item Lagerung: Extremität hängen lassen, weich und warm lagern
        \item Schonend transportieren
    \end{itemize}
\end{gMassnahmenEnv}

\subsubsection{Tiefe Venenthrombose}

\gDefinition{Blutgerinnsel (Thrombus) verstopft Vene}

\paragraph*{Ursachen:}

\gMargin{1
 }

\gMargin{2
}

\begin{itemize}
    \item Venöse Insuffizienz (Krampfadern = Venenaussackungen aufgrund
    gestörten  Blutrückflusses, Venenklappeninsuffizienz), Rückstau
    \item Immobilisation
    \item Bettlägrigkeit
    \item Lange Resien ohne Bewegung: \index{Economy class
    Syndrom}\gT{Economy class Syndrom}
\end{itemize}

\gDias
%1
{\includegraphics[width=\marginparwidth]{./images/economyclass.jpg}
\captionof{figure}{Die Economyclass. Sorgsam geschlichtet verbringen Menschen
hier Stunden damit Thrombosen zu basteln.
\gSourcePicture{Gabriel} \gAuthNotesPicLegalOk{}{}}}
%2
{\includegraphics[width=\marginparwidth]{./images/lovenox1.jpg}
\captionof{figure}{Thromboseprophylaxe. \gZ{Niedermolekulares Heparin
(hier "`Lovenox\texttrademark") verhindert Thrombosen die z.B. durch lange
Immobilisation (Reisen, Bettlägrigkeit, Liegegips, \ldots) enstehen können. Die
Substanz wird unter die Haut ("`subkutan"') gespritzt.
\gSourcePicture{Gabriel} \gAuthNotesPicLegalOk{}{}}}}
%3
{leer}

\paragraph{Symptome}

\begin{itemize}
    \item Schmerzen
    \item Schwellung (einseitig, anderes Bein auch ansehen zum Vergleich!)
    \item Dunkel-bläuliche (livide ) Verfärbung
    \item Überwärmung    
\end{itemize}

\paragraph{Gefahr}

\begin{itemize}
    \item Abschie{\ss}en des Thrombus ${\rightarrow}$
    \gEs{Lungenembolie} (s.
    I.II{\guillemotright}\pageref{bkm:RefHeading37612625})
\end{itemize}

\begin{gMassnahmenEnv}
    
    \paragraph{Spezielle Ma{\ss}nahmen}
    
    \begin{itemize}
        \item Lagerung:  liegend, Bein hoch!
        \item Bewegungsverbot: Pat. nicht mehr gehen lassen  (Emboliegefahr!)
    \end{itemize}
    
\end{gMassnahmenEnv}

\section{Respiratorische / Pulmologische  Notfälle}
\subparagraph{Störungen der Atmung - Übersicht}

\begin{center}
 [Warning: Image ignored] [width=3.754cm,height=4.22cm]{AASdev20090228-img75.png}
\gTempPicMissing{Splash Pulmo Notfälle img75}
\end{center}

\begin{itemize}
    \item mechanische Verlegung / Aspiration    
    \item Asthma bronchiale  / COPD
    \item Lungenembolie
    \item Lungen\"odem
    \item Lungenentz\"undung (Pneumonie)
\end{itemize}

\subsection{Mechanische Verlegung}

Durch eine Verlegung der Atemwege ist die Ventilation behindert bzw. aufgehoben.

Spezialfall: \index{Aspiration}\gT{Aspiration} =
Fremdk\"orper/-säfte (Magensaft, Erbrochenes, Bolus) gelangt in die
unteren Atemwege

Besonder gefährdet hinsichtlich einer Atemwegsverlegung sind:

\begin{figure}[h]
 \includegraphics{./images/pharynx-schema.png}
\caption{Atemwegsverlegung (Schema).
\gAuthNotesPicLegalNotOk{pharynx-schema}{}}
\end{figure}

\begin{itemize}
    \item Kinder
    \item bewu{\ss}tseinseingeschränkte Personen
    \item Z.n. Bienenstich bei Allergiker (Kehldeckel\"odem)
\end{itemize}

\paragraph{Symptome}

\begin{itemize} 
    \item Dyspnoe, Husten
    \item in/exspiratorisches Pfeifen (Stridor)    
    \item Atemstillstand, in weiterer Folge Kreislaufstillstand
\end{itemize}

\begin{gMassnahmenEnv}
    \paragraph{Spezielle Maßnahmen}
    
    \begin{itemize}
        \item Atemwege freimachen
        \item Esmarch-Handgriff
        \item Schlag zwischen Schulterblätter
        \item Heimlich-Griff
        \item \gTxBeiVit \gTxMassVitBes 
        \begin{itemize}
            \item Lagerung: Oberk\"orper hoch
            \item $O_2$ 10~l, bei Bedarf bis zu 15~l
        \end{itemize}
    \end{itemize}
    
\end{gMassnahmenEnv}



\subsection{Asthma bronchiale}

\gDefinition{Chronische Erkrankung mit mehr oder weniger
häufigen Episoden von Anfällen mit massviver Atemnot oder
Atembehinderung durch Verengung der Bronchien.}

Im akuten Asthmaanfall kommt es zu einer Verengung der
Bronchien, die Ausatmung wird massiv erschwert und der
Patient hat das Gefühl nicht genug Luft zu bekommen.
Mitunter ist ein pfeifendes Atemgeräusch wahrnehmbar. 

Der Patient nimmt miest automatisch und unbewusst eine
Position ein die seine Atemhilfsmuskulatur unterstützt: Er
sitzt aufrecht und stützt sich oft nach vorne oder hinten
ab. 

Asthma ist bei den meisten Patienten oft schon bekannt und
kann bzw. muß erfragt werden. Viele Patienten haben auch
einen oder mehrere Sprays oder Inhalationsgeräte zur Dauer-
oder Akutbehandlung. \gAuthNotes{Gabe von
Beta-Mimmetikum-Spray muss noch diskutiert werden.} 

Besonders betroffen vom Asthma bronchiale aind Patienten
mit bekannten Allergien. Folglich kann alles, was bei dem
Patienten eine Allergie verursacht auch einen Asthmaanfall
auslösen (Pollen, Gräser, Katzen, Zigarettenrauch, \ldots). 

Ist bei Eintreffen des Rettungsdienstes noch keine
Besserung der Atemnot eingetreten ist diese Atemnot als
massiv und lebensbedrohend einzuschätzen. Es ist nicht
damit zu rechnen daß der Zustand von alleine wieder vergeht
und bedarf einer medikamentösen Behandlung durch einen
Notarzt. 

\paragraph{Status Asthamticus} Der "`Status"' ist eine
gefürchtete Komplikation des Asthmaanfalles bei der die
Atemnot nicht wieder vergeht. Es besteht die Gefahr daß der
ÜPatient nach einiger Zeit keine Kraft mehr zum atmen hat
(Erschöüfung). Der Status asthmaticus kann tötlich enden. 

Wiederholte Anfälle von Atemnot durch

\begin{center}
[Warning: Image ignored] [width=3.282cm,height=4.502cm]{AASdev20090228-img77.png}
\gTempPicMissing{Asthma img77}
\end{center}
\begin{itemize}
    \item {\mdseries verengte Bronchien}
    
    \item {\mdseries vermehrte Schleimbildung}
    
\end{itemize}

\paragraph{ Symptome}

\begin{itemize}
    \item Akute Atemnot mit trockenem Husten
    \item Exspiratorischer Stridor (Brummen, Pfeifen, Giemen)
    \item Einsatz der Atemhilfsmuskulatur
    \item Verlängerte Ausatmungsphase
\end{itemize}

\paragraph*{Komplikationen}

\begin{itemize}
    \item Zyanose
    \item Status asthmaticus
    \item Panik (${\rightarrow}$ Tachykardie)
\end{itemize}

\begin{figure}[h]
    \includegraphics{./images/copd-bronchus.png}
    \captionof{figure}{Bronchus bei COPD.
    \gAuthNotesPicLegalNotOk{copd-bronchus}{}}
\end{figure}

\begin{figure}[h]
    
    \begin{captionbeside}{\gAuthNotesPicLegalOk{}{
    }Patientin mit einem akuten Asthmaanfall, sitzend auf einer Treppe. Man beachte wie sie sich hinten mit den Armen aufstützt. Die Erstmaßnahmen bei 
    vital bedrohten Patienten wurden ergriffen: ·~Situationsrechte Lagerung, 
    ·~Sauerstoffgabe, ·~Reanimationsbereitschaft (Defibrillator eingeschalten, 
    sonstiges Zubehör nicht am Foto gezeigt sichtbar), ·~Notarztnachforderung, 
    ·~Vitalwerte erheben und Monitoring: Blutdruck wird gerade gemessen, 
    die Patientin ist am Monitor angeschlossen. \gSourcePicture{Gabriel}}
    \includegraphics[width=6.5cm]{./images/asthma-paramedic.jpg}
    \label{fig:asthma-paramedic}
    \end{captionbeside}
\end{figure}

\begin{center}
            [Warning: Image ignored] [width=6.378cm,height=4.784cm]{AASdev20090228-img79.jpg}
            \gTempPicMissing{Asthma img79}
        \end{center}

\begin{gMassnahmenEnv}
    
    \paragraph{Spezielle Maßnahmen beim akuten
    Asthmaanfall} ~
    
     Einschätzen der vitalen Bedrohung. 
    \gEs{\gTxBeiVit} \gTxMassVitBes
    
    \begin{itemize}
        \item Lagerung: Oberk\"orper hoch, evt. abstützen
        lassen (Atemhilfsmuskulatur unterstützen)
        \item $O_2$ 10\,l, bei Bedarf bis zu 15\,l
    \end{itemize}
    
    \gEs{weiters:}
    
    \begin{itemize}
        \item Lagerung: Oberk\"orper hoch, evt. abstützen
        lassen (Atemhilfsmuskulatur unterstützen)
        \item $O_2$ \footnote{Achtung: In seltenen Fällen
        kann es vorkommen dass bei hochdosierter Sauerstoffgabe der Patient zu
        atmen aufhört! Siehe Kap. COPD (\ref{topic:copd})}
        ca. 6\,l, bei Bedarf bis zu 15\,l
        \item Vitalwerte erheben \& Monitoring 
        \item Beruhigen, voratmen (Lippenbremse)
        \item Sprays, die Patient ev. bei sich hat d\"urfen nicht mehr genommen
        werden.
    \end{itemize}
    \gAuthNotes{Einnahme von verordneten Sprays analog zu
    Nitro?\par $O2$-Gabe?}
\end{gMassnahmenEnv}





\subsection{COPD - Chronic obstructing pulmonary
disease}\label{topic:copd}

Syn.: Chronisch-obstruktive Lungenerkrankung.


\begin{figure}[h]
    \begin{center}
\includegraphics{./images/copd-stufen.png}
\caption{Eine COPD
entsteht nicht pl\"otzlich: Ein COPD{\textquotesingle}ler hat eine
Karriere hinter sich. \gSourcePicture{Quelle unbekannt} \gAuthNotesPicLegalNotOk{copd-stufen.png}{}}
\end{center}
\end{figure}



\subparagraph*{Sammelbegriff f\"ur}

\begin{itemize}
    \item Lungenemphysem (\"Uberblähung der Lungenbläschen)
    \item chron. obstruktive Bronchitis
\end{itemize}

\marginpar{weiters:
\begin{itemize}
    \item vermehrter Auswurf
    \item Ausatmung (Exspiration ) erschwert
    \item {\quotedblbase}Raucherhusten{\textquotedblleft}
\end{itemize}}
\subparagraph*{Symptome}

\begin{itemize}
    \item Belastungsdyspnoe (bei schwererer COPD auch in Ruhe)
    
    \item Husten mit Auswurf
    
    \item Brummendes, giemendes Ausatemgeräusch
    
    \item Fa{\ss}thorax
    
    \item {\bfseries\color[rgb]{0.5019608,0.0,0.0}
    Fa{\ss}thorax erklären}
    
    \item ben\"otigt evt. Heimsauerstoff
    
\end{itemize}
[Warning: Image ignored] {AASdev20090228-img81.png}
\gTempPicMissing{img 81}

Oben links: Bei der COPD sind die kleinen Luftwege verschleimt und
verengt.

Unten links: Die Lungenbläschen (Alveolen) sind \"uberbläht, weil
die Luft nur erschwert wieder entweichen kann. 

\paragraph*{Probleme mit Sauerstoff bei COPD-Patienten}
Die $CO_2$ - Konzentration im Blut ist bei COPD-Patienten chronisch hoch,
der K\"orper gew\"ohnt sich daran. Der K\"orper regelt dann den
Atemantrieb direkt \"uber den $O_2$-Gehalt im Blut. 

Wird der $O_2$-Gehalt pl\"otzlich k\"unstlich erh\"oht fehlt der
\index{COPD!Atemantrieb}Atemantrieb und der Patient kann einen
\gEs{Atemstillstand} erleiden. 

\gFazit{Sauerstoff vorsichtig dosieren, anfänglich nur 2-3~l/ min. und
    Atmung überwachen.
    
    Wenn der Patient bereits Heimsauerstoff benutzt kann man bei
    Atemnot versuchen 1-2~l/min höher zu dosieren. Die Atmung muss
    dabei überwacht werden! 
    
    Heimsauerstoff soll weiter gegeben werden!}

\marginlabel{
Sauerstoff vorsichtig dosieren, anfänglich nur 2-3~l/ min. und
    Atmung überwachen.
    
    Wenn der Patient bereits Heimsauerstoff benutzt kann man bei
    Atemnot versuchen 1-2~l/min höher zu dosieren. Die Atmung muss
    dabei überwacht werden! 
    
    Heimsauerstoff soll weiter gegeben werden!
}


\begin{gMassnahmenEnv}
    \paragraph{Spezielle Maßnahmen}

    \begin{itemize}
        \item
        Wenn bei Eintreffen noch keine Besserung eingetreten ist Notarzt
        erwägen!
        \item $O_2$ \footnote{Sauerstoff vorsichtig dosieren, anfänglich nur 2-3~l/ min. und
    Atmung überwachen. Wenn der Patient bereits Heimsauerstoff benutzt kann man bei
    Atemnot versuchen 1-2~l/min höher zu dosieren. Die Atmung muss
    dabei überwacht werden! Heimsauerstoff soll weiter gegeben werden!}
        
        \item Oberkörper hoch
        \item voratmen, Lippenbremse
        \item eigener Spray darf nicht mehr genommen werden
        
    \end{itemize}
    \gAuthNotes{Spray bei COPD -- Regelung wie bei Nitro??}
\end{gMassnahmenEnv}



\subsection{Lungenembolie}\label{topic:lungenembolie}

\gDefinition{Einschwemmen eines Gerinnsels in die  Strombahn des Lungenkreislaufs}

\subparagraph*{Herkunft des Embolus}

Der Embolus der die Lungenarterien verstopft entsteht
normalerweise nicht in den Lungengefäßen selbst sondern er
wird angeschwemmt. Sehr oft löst sich ein Thrombus bei
einer tiefen Beinvenenthrombose und schwimmt im Blut bis in
die Lunge. Daher gelten die
Risikofaktoren der tiefen Beinvenenthrombose
(Immobilisation, lange Reisen, \ldots) auch bei der
Lungenembolie. 

\marginlabel{\begin{itemize}\item Tiefe
(Bein-)Venenthrombose\item Immobilisation (Gips, lange Reise)\item Herz: Bei Rhythmusstörungen\item Tumorpatienten
\end{itemize}}Auch bei Herzrhytmusstörungen wie einem
Vorhofflimmern kann ein Thrombus im Vorhof selber entstehen, 
sich losreißen und Lungengefäße verstopfen. Bei
Gerinnungstörungen oder auch Krebserkrankungen kann das
Gerinnungssystem verrückt spielen und Thromben bilden die
zu einem Embolus werden.

\begin{figure}[hb]
    \caption{[Warning: Image ignored] [width=8.157cm,height=8.898cm]{AASdev20090228-img82.png}}
\end{figure}

\paragraph{Symptome}

\begin{itemize}
    \item Atemnot
    \item Akuter Thoraxschmerz
    \begin{itemize}
        \item eher atemabhängige Schmerzen (Atemstopp)
    \end{itemize}
    \begin{itemize}
        \item Angst
        \item Atemfrequenz und Herzschlag ${\uparrow}$
    \end{itemize}
\end{itemize}

Eine Lungenembolie kann man leicht mit einem Herzinfarkt oder Angina
pectoris verwechseln. Dies hat aber bei der Erstversorgung kaum
Konsequenz, die Ma{\ss}nahmen sind die gleichen.

\paragraph{Typische Befunde}

\begin{itemize}
    \item Zeichen einer tiefen Beinvenenthombose (mögliche
    Quelle des Embolus): Bein: einseitg geschwollen,
    überwärmt, rot/rot-bläulich verfärbt, schmerzhaft
    \item evt. gestaute Halsvenen
    
\end{itemize}

\subparagraph{Typische Zeichen}

\begin{itemize}
    \item Immobilisation oder Operation innerhalb
    der letzten 4 Wochen (\gEs{Lange Reise}, Flug, Auto,
    Bahn, Bettlägrigkeit, Gips, {\dots}
    \item 
    \item \gZ{Schwangerschaft}
    \item \gZ{Raucher(in)}
    \item \gZ{Verh\"utungspille}
\end{itemize}

\paragraph{Wann ist die Diagnose \emph{"`Lungenebmolie"'}
 wahrscheinlich?} nach
\ref{Wells2000}

\begin{itemize}
    \item {\bfseries Klinischer Verdacht} basierend auf
    Erfahrungswerten
    \item {\bfseries Andere Diagnosen sind weniger
    wahrscheinlich}
    \item Herzfrequenz \textgreater 100\,\nicefrac{}{min}
    (Tachykardie)
    \item Immobilisation oder Operation \gZ{innerhalb
    der letzten 4 Wochen} (\gEs{Lange Reise}, Flug, Auto,
    Bahn, Bettlägrigkeit, Gips, {\dots}
    \item Bereits einmal durchgemachte tiefe
    Beinvenenthrombose oder Lungenembolie
    \item Bluthusten
    \item Kebserkrankung \gZ{innerhalb der letzten 6 Monate}
\end{itemize}

Die ersten beiden, fett gedruckten, Kriterien sind die
wichtigsten Kriterien. 

\paragraph{Ma{\ss}nahmen}

\begin{itemize}
    \item \gTxMassVitBes
    \begin{itemize}
        \item Lagerung: Oberkörper hoch
        \item $O_2$ 10~l, bei Bedarf bis zu 15~l
    \end{itemize}
    \begin{itemize}
        \item Striktes Bewegungsverbot
        \item Beengende Kleidung öffnen
    \end{itemize}
\end{itemize}



\subsection{Lungenödem}\label{topic:lungenoedem}

{\mdseries
Fl\"ussigkeitsstau in den Alveolen und Interstitium}

\paragraph{Ursachen}

\begin{itemize}
    \item Durch R\"uckstau: akute Linksherzinsuffizienz, Linksherzversagen
    
    \item z.B. bei MCI, zusäztliche Belastung bei vorgeschädigtem Herz
    (Infekt)
    
    \item Durch Entz\"undung: Inhalation toxischer Substanzen
    
\end{itemize}
\gCave{
Das gefährliche am Lungenödem sind dessen Ursachen!}

\subparagraph*{
Symptome:}

\begin{itemize}
    \item Atemnot
    
    \item Zyanose
    
    \item \index{Brodelnde Atemgeräusche}Brodelnde Atemgeräusche, oft
    ohne Hilfsmittel hörbar, im Extremfall schon von der Wohnungst\"ur
    aus.
    
    \begin{center}

{\mdseries
Brodelnde Atemgeräusche: Als ob jemand mit einem Strohhalm in ein Glas
Wasser bläst.}


\end{center}
\end{itemize}

\begin{gMassnahmenEnv}
    
    \paragraph*{
    Spezielle Ma{\ss}nahmen}
    
    \begin{itemize}
        \item {\color[rgb]{0.0,0.0,0.5019608}
        Erstma{\ss}nahmen bei vital bedrohten Patienten}
        
    \end{itemize}
    \begin{itemize}
        \item {\color[rgb]{0.0,0.0,0.5019608}
        Lagerung: Oberkörper hoch, \gE{wenn möglich beine
        herabhängen lassen}}
        
        \item {\color[rgb]{0.0,0.0,0.5019608}
        O\gSub{2} 10~l, bei Bedarf bis zu 15~l}
        
        \item {\color[rgb]{0.0,0.0,0.5019608}
        Notarzt}
        
        \item {\color[rgb]{0.0,0.0,0.5019608}
        Reanimationsbereitschaft}
        
    \end{itemize}
    \begin{itemize}
        \item {\color[rgb]{0.0,0.0,0.5019608}
        Vitalwerte erheben \& Monitoring}
        
        \item {\mdseries\color[rgb]{0.0,0.0,0.5019608}
        Striktes Bewegungsverbot}
        
        \item {\mdseries\color[rgb]{0.0,0.0,0.5019608}
        Beengende Kleidung öffnen}
        
        \item {\mdseries\color[rgb]{0.0,0.0,0.5019608}
        \gZ{evt. Gabe von Nitroglycerin-Spray durch einen
        Notfallsanitäter}}
        
    \end{itemize}
    
\end{gMassnahmenEnv}

Patienten mit einem Lungenödem (v.a. bei kardialer Ursache) sind
\gEs{lebensbedrohlich krank}. Das Herz vebraucht oft
schon die allerletzten Reserven, schon die Aufregung durch den
Transport im Stiegenhaus zum Auto kann das Herz endg\"ultig in den
Herz-Kreislaufstillstand treiben. 

Es kommt häufig vor das diese Patienten \gEs{im
Stiegenhaus} oder \gEs{im Auto}
\gEs{reanimationspflichtig} werden. 

\gCave{Kein Transport ohne Notarzt! 

Auch wenn das Spital
{\quotedblbase}um{\textquotesingle}s Eck{\textquotedblleft}
ist. Der Patient schafft es möglicherweise nicht 'mal mehr
ins Auto!}



\subsection{Lungenentzündung --- \gT{Pneumonie}}

\begin{itemize}
    \item durch Keime: Bakterien, Pilze
    \gZ{(Pneumokokken,Chlamyden, Candida bei
    HIV/Immunsupression)}
    
    \item durch \index{Pneumonie!Aspiration}Aspiration
    
\end{itemize}

Höheres Lebensalter begünstigt Infektionen wie z.B. Pneumonien.

\paragraph*{
Symptome}

Die Symptome sind oft nicht charakteristisch, eine Diagnose kann meist
erst im Krankenhaus gestellt werden. 

\begin{itemize}
    \item 
    Dyspnoe, Husten
    
    \item 
    Fieber (mu{\ss} aber nicht sein)
    
    \item 
    Reduzierter Allgemeinzustand (AZ)
    
\end{itemize}

\begin{gMassnahmenEnv}
    
    \paragraph*{
    Spezielle Ma{\ss}nahmen}
    
    \begin{itemize}
        \item
        Symptomatisch
        
        \item
        Wärmeerhalt

        \item 
        Hospitalisierung, Abteilung: Innere Medizin
        
    \end{itemize}
    
\end{gMassnahmenEnv}


\subsection{Hyperventilationstetanie}\label{topic:hyperventilationstetanie}
\begin{itemize}
\item 90\% psychogen! (typisch: pubertierende  Mädchen)
\item Lungenembolie
\item diabetische Stoffwechselentgleisung (Ketoazidose, kompensatorisch)
\end{itemize}


\subparagraph{Symptome}

\begin{itemize}
\item Unruhe, Aufregung, evt. Angstgefühle, Panik
\item Tachypnoe mit Alkalose $\rightarrow$ durch Änderung
des pH-Werts Änderung des Kalziumspiegels im Blut, dadurch
\begin{itemize}
    \item Kribbeln, 
\item Krämpfe (Gesichtsmuskulatur),  Pfötchenstellung
\end{itemize}
\end{itemize}

\subparagraph{Therapie}

\begin{itemize}
\item {\mdseries
beruhigen!!!}

\item {\mdseries
Kommandoatmung}

\item {\mdseries
CO\gSub{2 }{}-R\"uckatmung  (Plastiksackerl,  Handschuh)}

\item {\mdseries
falls nicht erfolgreich, Notarzt f\"ur Sedierung  nachfordern}

\item {\mdseries
KEIN O{2}  verabreichen!!! Kontraindikation!!!}

\end{itemize}







\section{Neurologische Notfälle}
Erkrankungen, Notfälle und sonstige Störungen die das Nervensystem betreffen. Neurologische Störungen gehen oft mit einer Störung des Bewusstseins
oder Schmerzen einher.

Eine Übersicht und willkürliche Einteilung:

\begin{description}
    \item [Schlaganfall]
    \item [Vergiftungen] können fast alles machen, und neurologische Symptome erst recht
    \item [Schädel-Hirn-Trauma] Abkz. 'SHT'; 
    \item [Stoffwechselstörungen] welche Einfluß auf das Nervensystem nehmen, z.B. Zuckerentgleisungen 
    \item [Infektionen] können das Nervensystem beeinflussen, z.B. Verwirrtheit bei Fieber
    \item [Alter] beeinflusst auch das Nervensystem, im Rahmen einer Demenz kommt es z.B. zum Abbau der Denkleistung
    \item [Krampfanfälle] können vom zentralen Nervensystem ausgehen
    \item [Erkrankeungen des St\"utzapparates] Aufgrund der räumlichen Nähe zum Stütz- und Bewegungsapparates kann das Nervensystem in Mitleidenschaft gezogen werden, z.B. bei Bandscheibenvorfällen
\end{description}
\gAuthNotes{Es fehlt: Lumbago, BWS/HWS}

\subsection{Schlaganfall --- Apoplektischer Insult --- Hirninfarkt}
\index{Schlaganfall}\index{Insult}

\emph{Syn.: \index{Apoplektischer Insult}\index{Apoplektischer
Insult}Apoplektischer Insult, \index{Apoplexie}Apoplexie}

Grundsätzlich gibt es 2 Formen:
\begin{description}

    \item[>>Trockener<< Insult] Verschlu{\ss} eines Hirngefä{\ss}es mit anschlie{\ss}endem
    Absterben der versorgten Hirnregion ({\quotedblbase}trockener
    Insult{\textquotedblleft}), oder
    \item[>>Blutiger<< Insult] Spontane Hirnblutung, z.B. durch Platzen eines Hirngefäßes, kann spontan oder aufgrund eines Traumas auftreten.
    \item 
    Sonderform \index{TIA}TIA: \index{Transitorisch -- Ischämische
    Attacke}Transitorisch -- Ischämische Attacke
    \begin{itemize}
        \item Symptome wie bei Insult, vergehen aber innerhalb von 24 Stunden.
    \end{itemize}
\end{description}

\begin{figure}[h]
 [Warning: Image ignored] {AASdev20090228-img83.png} 
 \gTempPicMissing{Insult img83}   
\end{figure}
\begin{figure}[h]
 [Warning: Image ignored] {AASdev20090228-img84.png}
 \gTempPicMissing{Insult img84}
\end{figure}
\begin{figure}[h]
    [Warning: Image ignored] [width=9.526cm,height=5.428cm]{AASdev20090228-img85.png}
    \gTempPicMissing{Insult img85}
\end{figure}

\paragraph{Symptome}

Die Symptome sind davon bhängig welche Hirnregion betroffen ist und 
können sehr unterschiedlich sein.

\gE{'Klassische'} Symptome sind:

\begin{description}
    \item[Störungen der Motorik] sind sehr häufig
    \begin{itemize}
        \item Halbseitenschwäche Aufgrund des Ausfalls eines Zentrums im Gehirn kann es zu einer Kraftminderung einer Seite kommen, dies kann bis zu einer schlaffen, kompletten Lähmung führen.
        \item Herabhängender Mundwinkel als zeichen einer Lähmung der Gesichtsmuskulatur
        \item Schluckstörungen
    \end{itemize}
    \item[Störungen der Sprache] sind ebenfalls sehr häufig
    \begin{itemize}
        \item Verwschene Sprache
        \item {\quotedblbase}wirre Sprache{\textquotedblleft} Worte fehlen, Silben
        werden vertauscht, ...
        
    \end{itemize}
    \item[Störungen der Wahrnehmung]
    \begin{itemize}
        \item Sehstörungen, plötzliche Erblindung, Doppelbilder,
        Gesichtsfeldausfälle
    \end{itemize}
    \item[Störungen des Bewußtseins] von Somnolenz bis zur Bewusstlosigkeit
    
    \item[Allgemeine Störungen]
    \begin{itemize}
        \item Schwindel
        \item Plötzlicher, donnerschlagartiger Kopfschmerz. Klassisch bei spontaner Hirnblutung.
    \end{itemize}
\end{description}

\subparagraph*{
Differentialdiagnose}

\begin{itemize}
\item Hypoglykämie ${\downarrow}$, ${\uparrow}$ (${\rightarrow}$ Bei
neurologischen Symptomen immer messen!)
\item \gZ{Tumor}
\end{itemize}

\begin{gMassnahmenEnv}
    \paragraph{Ma{\ss}nahmen}
    
    \begin{itemize}
        \item \gTxMassVitBes 
        \begin{itemize}
            \item Lagerung: \gE{leicht erhöhter Oberkörper}
            \item $O_2$ 10~l, bei Bedarf bis zu 15~l
        \end{itemize}
        \item \gEs{Neurocheck} inkl. \gEs{BZ-Messung} (DD!)
        \item Hospitalisation: \gEs{Stroke Unit!}
    \end{itemize}
    
\end{gMassnahmenEnv}

\subsection{Bakterielle \index{Meningitis}Meningitis}

\gDefinition{Eine eitrige Entz\"undung der Hirnhäute aufgrund einer Entzündung
mit Bakterien.}

Die bakterielle Meningitis ist eine gefürchtete Erkrankung. Einerseits ist sie
hoch ansteckend, andererseits kann sie rasch tödlich enden. Gefährdet sind
grundsätzlich alle Altersgruppen, besonders aber Kinder, Jugendliche und alte
Menschen.
 
\begin{wrapfigure}{r}{5cm}

\includegraphics[width=\linewidth]{./images/meningitis-kind-waterhouse.png}
\captionof{figure}{Endstadium einer bakteriellen Meningitis.
Die Flecken auf der Haut sind Blutungen in Folge einer nicht mehr
funktionierenden Blutgerinnung. Der Tod ist bei diesem Patienten
innerhalb weniger Stunden zu erwarten.\newline
\gSourcePicture{unbekannt}}\label{fig:meningitis-baby}
\end{wrapfigure}
Typische Symptome sind Nackensteifigkeit, Fieber und Lichtscheue. \gZ{Im
fortgeschrittenen Stadium kann es zu massiven Gerinnungsstörungen kommen, am
Körper bilden sich dann dunkle Hämatome (s. Abb. \ref{fig:meningitis-baby}.
Einer der häufigsten Erreger sind die Meningokokken, man spricht dann von einer Meningokokken-Meningitis.}


\paragraph{Symptome}

\begin{itemize}
    \item Kopfschmerzen, Lichtscheue
    \item Erbrechen, \"Ubelkeit
    \item Nackensteifigkeit (Meningismus)
    \item evt. Ausschlag (Exanthem )
\end{itemize}


\gAuthNotes{FEHLT: Querverweis Meningimus\\FEHLT: Ma{\ss}nahmen}

\subsection{Vom Gehirn ausgehende Krämpfe --- Zerebrale Krampfanfälle}

\gDefinition{Vom Gehirn ausgehende motorische Krampfanfälle aufgrund von akuten
oder chronischen Erkrankungen des Gehirns. Dabei entladen sich unkontrolliert Nervenzell-Gruppen im Gehirn.}

Die zerebralen Krampfanfälle kann man auf verschiedene Arten einteilen

\begin{enumerate}
    \item nach Art:
    \begin{description}
        \item[tonisch]
        \item[klonisch]
    \end{description}
    \item nach Lokalisation:
    
    \begin{description}
        \item [Fokaler Anfall] Der Krampfanfall ist auf eine Körperregion beschränkt, z.B. auf eine Hand. Der Patient kann bei Bewußtsein sein.
       Ein fokaler Anfall kann sich ausweiten und in einen generalisierten Anfall übergehen. \gZ{Der fokale Anfall wird auch \emph{"`petit mal"'}
      (\emph{franz.:} "`kleines Unglück"') genannt.}
        
        \item [Generalisierter Anfall] Er betrifft den ganzen Körper. Der
        Patient ist nicht bei Bewußtsein, es besteht während des Anfalls
        ein Atemstillstand. Der Patient kann starke Kräfte entwickeln es besteht eine massive Verletzungsgefahr. \gZ{Der generalisierte Anfall wird
       auch \emph{"`grand mal"'} (\emph{franz.:} "`großes Unglück"') genannt.}   
        
    \end{description}
    \end{enumerate}

\paragraph{Ursachen}\marginlabel{Ursachen}

Für zerebrale Anfälle gibt es eine Reihe von unterschiedlichen Ursachen die
direkt oder indirekt das Gehirn betreffen. Einige betreffen die Hirn- und
Nervenstruktur, andere sorgen dafür dass das Hirn zuwenig Sauerstoff und andere
Nährstoffe erhält. 

\begin{description}
    \item[Epilepsie]Die Epilepsie ist eine chronische
    Erkrankung welche durch wiederkehrende Episoden von zerebralen Krampfanfällen gekennzeichnet ist.
    \item[Vergiftungen](Medikamente/Drogen/...) können wieder fast alles machen, also auch Krampfanfälle.
    \item[Hirntumor] Durch den Tumor wird das Hirn beschädigt, auch dadurch kann es zu Anfällen kommen.
    \item[Narbengewebe] nach Insult oder Hirnverletzung.
    \item[Hirnblutung]
    \item[Sauerstoffmangel] z.B. Ischämie aufgrund eines Gefäßverschlusses
    \item[BZ] Der Dauerbrenner. Eine Blutzuckerentgleisung kann fast alle neurologischen Symptome vortäuschen!
    \item[angeboren] ohne organisch fa{\ss}bare  Ursache
    \item[Fieberkrampf] Fieber bei Kleinkindern
    \item[Infektionen]
\end{description}

\paragraph{Fokaler Anfall}
Der Krampfanfall ist auf eine Körperregion beschränkt, z.B. auf eine Hand. Der Patient kann bei Bewußtsein sein.
       Ein fokaler Anfall kann sich ausweiten und in einen generalisierten Anfall übergehen. \gZ{Der fokale Anfall wird auch \emph{"`petit mal"'}
      (\emph{franz.:} "`kleines Unglück"') genannt.}
      


\paragraph{Generalisierter Anfall}

Er betrifft den ganzen Körper. Dem Anfall kann ein
\gE{Dämmerzustand}\marginlabel{Dämmerzustand -- Absence}
vorausgehen, \gZ{eine sog. "`Absence"'}. Der Patient wirkt
dabei plötzlich abwesend und \emph{"`irgendwie komisch"'}.
Dann verliert er das Bewußtsein und die Muskeln werden
anfangs schlaff. Wenn der Patient gestanden ist fällt er
ungebremst auf den Boden (Cave: \gEs{Verletzungen!}).
Manchmal kommt es vor daß der Patient einen
Initialschrei\index{Initialschrei} von sich gibt, der
gruselig klingen kann. Dann kommt es zu
\gEs{tonisch-klonischen Krämpfen}\marginline{Tonisch-klonische Krämpfe}
die den ganzen Körper betreffen. Während des Anfalls ist der
Patient nicht bei Bewußtsein, es besteht ausserdem ein
Atemstillstand. Er kann jedoch starke Kräfte entwickeln, es
besteht massive
Verletzungsgefahr\marginline{Verletzungsgefahr}.

Der Krampf dauert meist wenige Minuten und hört von selbst
auf. Der Patient ist danach erstmal bewusstlos, im weiteren
Verlauf mit mehr oder weniger großem Aufwand erweckbar. Er
beindet sich dann in der sogenannten
\gT{Nachschlafphase}\index{Nachschlafphase}\marginline{Nachschlafphase}.
Der Patient \gE{klart im Verlauf mehr und mehr auf}. Es ist
wichtig zu versuchen mit dem Patienten in Kontakt zu treten
um seinen Wachheitsgrad \marginline{Aufklaren}und dessen
Verlauf beurteilen zu können. Für den Anfall selbst besteht
Amnesie! Daher muß man den Patienten aufklären was passiert
ist, wo er sich befindet, etc. Ein \gEs{Neuro-} und
\gEs{Traumacheck} ist nach \underline{\gE{jedem}} Anfall
unbedingt notwendig!

\gZ{Der generalisierte Anfall wird auch \emph{"`grand mal"'} (\emph{franz.:} "`großes Unglück"') genannt.}

\paragraph{Fieberkrampf}\label{topic:fieberkrampf}

Krampfanfall bei Kleinkindern bei schnellem, hohen Fieberanstieg
({\textgreater} 39{\textdegree}C)

Sieht aus wie generalisierter grand mal-Anfall, NUR bei
Kindern {\textless} 4 Jahre.

\paragraph{Status Epilepticus} Der Status Epilepticus ist eine
lebensgefährliche\marginline{\gEs{LEBENSGEFAHR!}} Steigerung des
generalisierten Krampfanfalls. Er besteht dann wenn der Anfall nicht nach kurzer Zeit vergeht oder innerhalb kurzer Abstände
mehrere Anfälle auftreten (z.B. 2 Anfälle innerhalb von 24 Stunden). Der Krampf
muß dann mit Medikamenten durchbrochen werden. 

\gEs{Im Rettungsdienst gilt praktisch jeder Krampfanfall den das Personal
selber miterlebt als Status-Verdächtig.}

\begin{gMassnahmenEnv}
\paragraph{Spezielle Ma{\ss}nahmen beim krampfenden Patienten}
\begin{itemize}
    \item \gTxMassVitBes
    \begin{itemize}
        \item Lagerung: sichern
        \begin{itemize}
            \item Patient vor Verletzungen sch\"utzen
            \item Möbel etc. entfernen oder polstern
            \item NICHT festhalten -- aber sichern, Eigenschutz
            beachten!\footnote{\textbf{Sichern}: Patient darf nicht von der
            Trage oder aus dem Bett fallen, der Kopf darf nicht am Boden oder Tischbein anschlagen, etc.}
        \end{itemize}
    \end{itemize}
    \item nach Anfall BAK-Schema, Atemwege freihalten
    \item stabile Seitenlage  (Erholungsschlaf wie  Bewu{\ss}tlosigkeit)
    \item Besondere Diagnostik:
    \begin{itemize}
        \item Neurocheck inkl. \gEs{BZ}!
        \item Traumacheck
    \end{itemize}
    \item Reizabgeschirmter Transport
    \item Hospitalisierung: Abt.f. Innere Medizin. (evt. auch
    Intensivüberwachung)
\end{itemize}
\end{gMassnahmenEnv}

\begin{gMassnahmenEnv}
\paragraph{Spezielle Ma{\ss}nahmen beim \emph{nicht-}krampfenden Patienten}
\begin{itemize}
    \item Umstände klären und entscheiden ob \gEs{Notarzt} notendig ist. NA wenn:
    \begin{itemize}
        \item 2 oder mehr Krampfanfälle innerhalb von 24 Stunden
        \item Anfälle "`erheblich anders"' als sonst üblich
        \item Erster Krampfanfall überhaupt
    \end{itemize}
    \item Besondere Diagnostik:
    \begin{itemize}
        \item Neurocheck inkl. \gEs{BZ}!
        \item Traumacheck
    \end{itemize}
    \item Reizabgeschirmter Transport
    \item Hospitalisierung: Abt.f. Innere Medizin.
\end{itemize}
\end{gMassnahmenEnv}




\section{Metabolische Störungen (Stoffwechselstörungen)}

Die relevanteste Stoffwechselstörung im Rettungsdienst ist der
\index{Diabetes mellitus}\gT{Diabetes mellitus}. Dieser
stellt eine Störung des Kohlehydrat-(zucker-)stoffwechsels dar.
\gEs{Zucker ist ein Kohlehydrat}. 

\subsection{Diabetes Mellitus}\label{topic:diabetes}

\gEs{Merke}: Diabetes mellitus ist eine chronische
Stoffwechselerkrankungen mit  gestörter Blutzuckerregulation. 
Durch die Erkrankung oder die unsachgemäße Behandlung kann es zu 
folgenden Notfällen kommen: 

\begin{description}
\item[\gT{Hyp}\gT{o}\gT{glykämie}] ... Unterzucker

\item[\gT{Hyp}\gT{er}\gT{glykämie}] ... {\quotedblbase}\"Uberzucker{\textquotedblleft}

\end{description}

\paragraph{Insulin}

Insulin ist ein Hormon und wird in den \gT{Beta-Zellen}
der Bauchspeicheldr\"use (Pankreas) produziert. Es sorgt daf\"ur dass
Zucker (\gT{Glukose}) aus dem Blut in die Zellen
aufgenommen wird. 

Zwei Störungen sind möglich:

\begin{enumerate}
    \item Es wird zu wenig Insulin vorhanden.
    
    \item Der Körper hat einen gesteigerten Bedarf an Insulin weil er schon 
    zu sehr daran gewohnt ist. Der Körper reagiert auf die 'normale' Dosis
    nicht mehr so wie er sollte, er braucht grö{\ss}ere Mengen Insulin um
    den gleichen Effekt wie ein Gesunder zu erreichen (Gewöhnungseffekt),
    irgendwann kannn aber der Körper nicht mehr Insulin produzieren.
    
\end{enumerate}

Folglich unterscheidet man 2 Typen:

\begin{description}
    \item[Diabetes mellitus Typ I]  \gZ{(IDDM,
    insuline dependent)} ... keine Insulinproduktion des Körpers, der Patient muss
    immer Insulin als Ersatz spritzen.
    Störung des Immunsystems ${\rightarrow}$ Zerstörung der
    Insulin-produzierenden ${\beta}$-Zellen des Pankreas ${\rightarrow}$
    kein Insulin -- absoluter Mangel
    
    \item[\gT{Diabetes mellitus Typ II}]
    \gZ{(NIDDM, non insuline dependent)} ... Zuviel Zucker
    (schlechte Ernährung), folglich zu viel Insulin, Körper gewöhnt
    sich, s.o. ${\rightarrow}$ braucht erst Sport, Medikamente, im
    fort\-ge\-schrittenen Stadium oft Insulin
    Relativer Insulinmangel, Körper ist \"ubersättig an Insulin und
    reagiert kaum mehr daruf \gZ{(Insulinresistenz)},
    später ${\beta}$-Zellen erschöpft
    
    \begin{itemize}
        \item späte Manifestation: {\textgreater}35 Jahre
        \item \"Ubergewicht
        \item 95\% aller Diabetiker
        \item 25-30\% erblich
    \end{itemize}
    
\end{description}

\paragraph*{
Spätfolgen:}

Ein hoher Blutzuckerspiegel \"uber Jahre schädigt den gesammten
Körper. Es kommt zu:

\begin{itemize}
    \item {\mdseries
    Gefä{\ss}schäden durch Verkalkung}
    
    \begin{itemize}
        \item Gefä{\ss}verschl\"usse (kleine und gro{\ss}e Gefä{\ss}e)
        \item \gEs{Beine}: {\quotedblbase}Diabetischer
        Fu{\ss}{\textquotedblleft}, in Folge der Mangeldurchblutung in der Peripherie heilen auch kleine
        Wunden nicht mehr, in weiterer Folge muss immer weiter amputiert
        werden.
        \item \gEs{Herz}: Die Herzkranzgefä{\ss}e
        (Koronar\-arterien) sind ebenso betroffen. Diabetiker sind Risikopatienten f\"ur
        Herz-Kreislauf-Erkrankungen wie zB Herz\-infarkt.
        \item \gEs{Hirn}: Wie beim Herzen kommt es zu
        Gefä{\ss}schäden, daher ist das Risiko f\"ur Schlagänfälle erhöht.
        \item \gEs{Niere}: Die sehr kleinen Gefä{\ss}e der
        Niere werde geschädigt, im weiteren Verlauf stellt die Niere ihre Arbeit ein
        ${\rightarrow}$ {\quotedblbase}Chronische
        Niereninsuffizienz{\textquotedblleft}. Diese Erkrankung macht einen
        Gro{\ss}teil der Krankentransporte aus.
    \end{itemize}
    \item {\mdseries
    Nervenschäden \gZ{(periphere Polyneuropathie )}:
    Gef\"uhlslosigkeit vor allem in den Extremitäten, aber auch am
    gesammten Körper: Ein Herzinfarkt kann zum Beispiel schmerzlos
    ablaufen!}
    
    \item {\mdseries
    Augenschäden -- Erblindung}
    
\end{itemize}

Krankheiten die sonst schmerzhaft sind können beim Diabetiker aufgrund
der Nervenschäden ohne Schmerzen ablaufen. 

\gCa{Achtung}: Stummer Herzinfarkt beim Diabetiker!



\subparagraph[Diabetes \ mellitus I]{Diabetes  mellitus I}

Manifestationsalter: 0-25

5\% der Diabetiker, normalgewichtig ({\quotedblbase}jung,
sportlich{\textquotedblleft}), 2-6\% erblich

\subparagraph{Therapie}

{\mdseries
Insulin (mu{\ss} gespritzt werden!)}

{\mdseries
Vorsicht: Hypoglykämie-gefährdet (wenn zuviel Insulin gespritzt oder
zu wenig gegessen wurde) ${\rightarrow}$  Spritz/E{\ss}abstand
er\-fra\-gen}

\subparagraph[Diabetes \ mellitus II]{Diabetes  mellitus II}

\subsubsection{Hypoglykämie}

\paragraph{Ursachen}

\begin{center}
[Warning: Image ignored] % Unhandled or unsupported graphics:
%\includegraphics[width=4.234cm,height=7.526cm]{AASdev20090228-img88.png}
\gTempPicMissing{DM img88}

\end{center}
\begin{itemize}
    \item {\mdseries
    zuviel Insulin (spritzen ohne zu Essen)}
    
    \item {\mdseries
    fasten (spritzen ohne zu Essen)}
    
    \item {\mdseries
    Sport}
    
    \item {\mdseries
    Alkohol}
    
\end{itemize}

Der Blutzucker ist normalerweise 80-110~mg/dl

\paragraph*{Symptome}

\begin{itemize}
    \item {\mdseries
    Unruhe, Zittern, Schwei{\ss}ausbruch}
    
    \item {\mdseries
    Hei{\ss}hunger, \"Ubelkeit}
    
    \item {\mdseries
    Tachykardie, Tachypnoe}
    
    \item {\mdseries
    Haut feucht/bla{\ss}}
    
    \item {\mdseries
    Verwirrtheit, Krampfanfälle, Bewu{\ss}tlosigkeit}
    
\end{itemize}

Hypoglykämie kann alle neurologischen  Störungen
aus\-lösen/imitieren! Bei jeder  neurologischen Symptomatik  (v.a. Bewu{\ss}tseinsstörungen)
an  Hypoglykämie denken!

\gFazit{${\rightarrow}$ BZ-Messung bei jeder neurologischen Störung. }



\begin{gMassnahmenEnv}
    
    \paragraph{Spezielle Maßnahmen}

Es ist wichtig zu erheben ob der Patient vital bedroht ist oder nicht. 
Eine vitale Bedrohung ist zum Beispiel bei Bewusstlosigkeit gegeben.

Stark getrübten Patienten soll nichts zu essen gegeben werden da Aspirationsgefahr besteht. 

\begin{description}
    \item [Evt. Standardmaßnahmen bei vital bedrohten Patienten]
    
    \item [BZ messen!] Zur Diagnosestellung \gE{vorher} und und \gE{nach} Zuckergabe
    \item [Zucker zuführen] Wenn der Patient bewusstseinsklar ist Traubenzucker geben. Danach nachessen
    lassen da der Traubenzucker nur eine kurze Wirkdauer hat.
    
    \item [evt. Notarzt] bei vitaler Bedrohung oder wenn Zuckergabe nur über eine venöse Gabe erfolgen kann
    (getrübte Patienten die aufgrund der Aspirationsgefahr
    nichts mehr essen dürfen).
    
\end{description}
\end{gMassnahmenEnv}

\subsubsection{Hyperglykämie}

\subparagraph*{Diabetisches Koma }

{\mdseries
(BZ {\textgreater} 400 mg/dl)}

ketoazidotisches Koma

\begin{itemize}
\item BZ ${\uparrow}$

\item saure Ketonkörper ={\textgreater} Azidose (\"Ubersäuerung)

\item kompensatorisch Ku{\ss}maul{\textquotesingle}sche Atmung
(CO\textsubscript{2} - Abatmung)

\marginlabel{\textbf{Erinnere Dich}: Hyperventilationstetanie ${\rightarrow}$
Alkalose (pH wird basisch). 

Bei der diab. Ketoazidose (sauer) versucht der Körper diesen
Mechanismus auszunutzen: 

{\quotedblbase}\textbf{Er will sich vom sauren pH zum basischen pH
atmen.{\textquotedblleft}}

${\rightarrow}$ Er hyperventiliert.
} %marginlabel

\item eher bei DM I

\end{itemize}

\paragraph{Hyperosmolares Koma}

\begin{itemize}
    \item BZ ${\uparrow}$${\uparrow}$
    
    \item zuviel Glucose ={\textgreater} rel. Fl\"ussigkeitsmangel
    
    \item Niere gestört
    
    \item DM II
    
\end{itemize}

\subparagraph*{Ursachen}

\begin{itemize}
\item Diätfehler (Kohlenstoff-reiche Ernährung)

\item Insulinzufuhr unzureichend

\item Stre{\ss}

\item Infektionen (Insulinbedarf ${\uparrow}$)

\end{itemize}

\paragraph{Symptome}

\begin{itemize}
    \item Durstgef\"uhl
    \item vermehrter Harndrang, Harnmenge ${\uparrow}$
    \item Bewu{\ss}tseinstr\"ubung bis Bewu{\ss}tlosigkeit
    \item Ku{\ss}maul{\textquotesingle}sche Atmung (sehr tief und schnell)
    \item Aceton-Geruch (Achtung, nicht mit Schnapsfahne verwechseln!)
    \item DM I: ausgeprägte Symptome
    \item DM II: schleichender Verlauf
\end{itemize}

\paragraph*{Therapie}

\begin{itemize}
\item Basistherapie Bewu{\ss}tlosigkeit

\item Notarzt

\item $O_2$

\end{itemize}

\clearpage
%%%%%%%%%%%%%%%%%%%%%%%%%%%%%%%%%%%%%%%%%%%%%%%%%%%%%%%%%%%
%%% Abdomen                                             %%%
%%%%%%%%%%%%%%%%%%%%%%%%%%%%%%%%%%%%%%%%%%%%%%%%%%%%%%%%%%%

\section{Abdominelle Notfälle}

\begin{wrapfigure}{r}{5cm}
    \includegraphics[width=\linewidth]{./images/operation.jpg}
    \captionof{figure}{\gAuthNotesPicLegalOk{}{}Oft die
    Abteilung der Wahl für einen abdominellen Notfall: Die
    Chirurgie. Aber auch die Abteilungen für
    Kinderheilkunde, Kinderchirurgie oder Gynäkologie (!) können manchmal
    gute Ziele sein. 
    \gSourcePicture{Gabriel}}
\end{wrapfigure}Der Bauch ist eine echte Herausforderung. Er
beinhaltet viele unterschiedliche Organe und Strukturen die
weh tun und
erkranken können. Viele dieser, oft auch mit starken Schmerzen verbundenen, Störungen
sind nicht übermäßig schlimm, einige sind jedoch sehr gefährlich oder können rasch
gefährlich werden. 

\paragraph{Anamnese und Untersuchung} Es ist daher wichtig
typische Symptome zu erheben und Alarmzeichen sofort zu erkennen.

Wesentlich ist bei allen unklaren Bauchbeschwerden die Erhebung einer exakten Stuhlanamnese:

\begin{itemize}
    \item Auffälligkeiten? (Geruch, Farbe, Konsistenz, Frequenz)
    \item Wann erfolgte der letzte Stuhlgang?
    \item Winde?
    \item Genaue Beschreibung von Harn und eventuell Erbrochenem.
    \item Bei Frauen im gebärfähigem Alter Regelanamnese.
\end{itemize}

Der Patient muß bis zur endgültigen Klärung der Situation
\gEs{nüchtern} belassen werden!

Wenn es der Patient toleriert muß eine Betrachtung und
Abtastung des unbekleideten Abdomens beim liegenden
Patienten erfolgen wie unter
\ref{topic:untersuchung-abdomen-abtasten} beschrieben.

\subsection{Der schmerzende Bauch}

Gleich vorweg: In vielen
Fällen wird der Grund des Bauchschmerzes nicht mit
hinreichender Sicherheit feststellbar sein und es bietet sich
keine sichere Verdachtsdiagnose an. Statt einer Diagnose wird
dann das Beschwerdebild angegeben (Ort des Schmerzes, Art des
Schmerzes, Dauer, Verlauf, Austrahlung, Ergebnis der
abdominellen Tastuntersuchung: Abdomen weich oder hart,
Druckschmerz, Loslassschmerz).

\minisec{Beispiele:}

\begin{quote}
{\sffamily\bfseries Diagnose:} {\ttfamily Unklarer
Bauschmerz: Kolikartige Schmerzen im linken Unterbauch seit 1/2h }

{\sffamily\bfseries Anmerkungen:} {\ttfamily keine Ausstrahlung, Druckschmerz im li. Unterbauch, Abdomen sonst  weich.}
\vspace{0.5em}

{\sffamily\bfseries Diagnose:} {\ttfamily Unklarer
Bauchschmerz: Brennend-stechender Schmerz im mittleren
Oberbauch seit 2h mit Ausstrahlung in den Rücken}

{\sffamily\bfseries Anmerkungen:} {\ttfamily Abdomen weich,
keine Druck- oder Loslassschmerzen“}
\vspace{0.5em}

{\sffamily\bfseries Diagnose:} {\ttfamily Unklarer
Bauchschmerz: „Dumpfer“ Schmerz im linken und rechten
Oberbauch seit gestern. }

{\sffamily\bfseries Anmerkungen:} {\ttfamily Abdomen
weich, keine Druck- oder Loslassschmerzen}
\end{quote}

Dennoch muß man bei jeden Patienten versuchen so gut als
möglich die möglichen Krankheitsbilder einzugrenzen und
gefährliche Erkrankungen nach Möglichkeit auszuschließen. Im
folgenden sind einige häufige Krankheitsbilder beschrieben,
die \gE{erkennbar}  und relativ häufig sind.

\subsection{Häufige und gut erkennbare Erkrankungen}

\subsubsection{Das Akute Abdomen --- Alarmzeichen mit Bretthartem
Bauch}\label{topic:akutes-abdomen}



\gDefinition{\gT{Akutes Abdomen}: Zustand mit plötzlichen,
starken \gEs{Bauchschmerzen} und \gEs{brettharten Bauch}, es
sind viele Ursachen  möglich!}

\gZ{Die Definition des Akuten Abdomens ist oft nicht ganz
eindeutig: Einerseits kann man da\-runter jede plötzlich auftretende
Baucherkrankung verstehen\footnote{Longmore et.al.: Oxford Handbook of
Clinical Medicine, 7\textsuperscript{th} edt., Oxford University Press,
2007\cite{Longmore2007}}\footnote{\cite{Renz-Polster2006}
S.611}.} 

Allerdings ver\-stehen viele Leute darunter eher
eine bedrohliche Bauch\-erkrankung, bei der man akut (schnell) handeln muss (Ab\-wehrspannung, harter Bauch). In der
Notfallmeidizin halten wir uns an letztere Definition. 


\gMarried
{Das „Akute Abdomen“ ist ein interventionspflichtiger
Symptomenkomplex, der auf einem sich akut oder chronisch
entwickelnden (Stunden, Tage oder auch Wochen) pathologischen
Prozeß im Bauchraum beruht.

Es hat grundverschiedene Ursachen, die eine gemeinsame
Symptomatik haben.

Der Patient liegt meistens im Bett. Oft ist
eine Schonhaltung (angezogene Beine, gekrümmte
Körperhaltung) zu beobachten. Er klagt über Schmerzen im
Bauch die können als diffus (im gesamten Bauch) oder auf ein Bauchsegment
beschränkt, angegeben werden. 

Die harte Bauchdecke („bretthart“) ist neben den Schmerzen
das Leitsymptom des Akuten Abdomens. 

Es besteht oft Übelkeit oder auch Erbrechen. Der Kreislauf
kann kompensiert sein, oder Zeichen von Schwäche zeigen,
bishin zu einer Schocksymptomatik (Blutverlust, Sepsis).} 
{
\gTabHeading{Symptome}

\begin{itemize}
    \item \gEs{Schmerz}: diffus (ganzer Bauch) oder
    lokalisierbar
    \item \gEs{Harte  Bauchdecke} (bretthart)
    \item Schonhaltung  (angezogene Beine, Patient
    kr\"ummt sich)
    \item \"Ubelkeit, Erbrechen
    \item evt. Stuhl-, Windverhalten
    \item evt. Fieber
    \item evt. Schock
\end{itemize}
}


Diese Auflistung der Symptome ist in der Praxis nicht
unbedigt vollständig bei einem Patienten nachzuweisen. Die
Intensität und die Anzahl der angeführten Symptome ist
einerseits vom Allgemeinzustand, vom Lebensalter, von
anderen vorbestehenden Erkrankungen des Patienten und der
Dauer des aktuellen Zustandes abhängig.

Wesentlich ist nach Eintreffen beim Patienten ist die
Beantwortung der Fragen \emph{"`Wie schlecht, bzw. wie gut
geht es dem Patienten? Ist der Patient akut vital
bedroht?"'}.



\paragraph*{Ursachen}

Grundsätzlich kommt fast jede Erkrankung oder Verletzung in
Betracht die Strukturen und Organe im Bauchraum betrifft,
entsprechend vielfältig sind auch die Ursachen: 

\begin{multicols}{2}
\begin{itemize}
    \item Bauchfellentz\"undung (Peritonitis)
    \item Darmverschlu{\ss} (Ileus)
    \item Blinddarmentz\"undung (Appendizitis)
    \item Entz\"undung der Bauchspeicheldr\"use  (Pankreatitis)
    \item Entz\"undung des Magens / Magengeschw\"ur
    \item Durchbruch / Zerreißung eines Hohlorganes
    \item Gallenkolik
    \item Blutungen in den Bauchraum
    \item Bauchaortenaneurysma
    \item Mesenterialinfarkt
    \item andere abdominelle Erkrankungen
    \item Entz\"undung der Leber (Hepatitis )
    \item Gynäkologische Erkankungen
    \item \ldots
\end{itemize}
\end{multicols}


\begin{gMassnahmenEnv}
    
\paragraph{Allgemeine Ma{\ss}nahmen (Ursachen-unabhängig)}

\begin{itemize}
    \item \gEs{Vitale Bedrohung einschätzen.} \gTxBeiVit
    \gTxMassVitBes
    \begin{itemize}
        \item $O_2$ 10\gLMin
        \item Lagerung: je nach Situation: bauchdeckenentspannend,
        Beine hoch, \ldots
    \end{itemize}
    \item Lagerung: grundsätzlich bauchdeckenentspannend
    (Knierolle, Beine angezogen), oder nach Wunsch des Patienten
    \item n\"uchtern lassen (in Hinblick auf möglicherweise bevorstehende
    Operation, nichts essen oder trinken lassen)
    \item Monitoring. Der Zustand kann sich abrupt
    verschlechtern.
    \item Hospitalisierung: ad Chirurgie, bei Kindern ${\rightarrow}$ Kinder, bei
    Frauen evt. Gynäkologie oder Chirurgie mit Gyn im Haus (SMZ-O, Rudi,
    Wili, etc.)
\end{itemize}

\end{gMassnahmenEnv}


\subsubsection{Darmverschluß -- Ileus}\index{Ileus}

\gDefinition{Unterbrechung der normalen  Darmpassage}

\gZ{\index{Ileus}\gT{Ileus}  (vollständig)

\index{Subileus}\gT{Subileus}  (unvollständig)}

\gMarried
{Typisch für den Darmverschluß ist die mangelhafte
Stuhlproduktion. Es kommt zu diffusen Bauchschmerzen und einem
Blähbauch, sowie zu Übelkeit und Erbrechen bishin zum
Erbrechen von zurückgestautem Kot (selten).

Der Verschluß kann mechanisch bedingt sein (durch
Hindernisse wie Verwachsungen oder Tumore), oder durch eine
Lähmung der Darmmuskulatur (Medikamente, Gifte, Entzündung
\ldots) verursacht sein. Bestimmte Schmerzmittel
\gZ{(Opiate)}, welche oft bei Krebspatienten eingesetzt
werden, sidn bekannt dafür Darmlähmungen zu erzeugen. }
{\gTabHeading{Arten}

\begin{itemize}
\item Mechanisches Hinderniss (Verschlingung, Abklemmung 
etc.)
\item Lähmung der Darmmuskulatur (Medikamente, Entzündung)
\end{itemize}

\gTabHeading{Symptome}

\begin{itemize}
\item Schmerzen
\item Stuhlverhalten
\item Erbrechen (im schlimmsten Fall Koterbrechen)
\end{itemize}}




\subsubsection{Blinddarmentzündung -- Appendizitis}



\gDefinition{Entz\"undung des \gEs{Wurmfortsatzes}
(\gEs{Appendix} \gZ{vermiformis}) des Blinddarms}

\gMarried {\gACh{GABG}{2009-05-07}{NEU}Im Vordergrund der akuten Appendizitis
steht anfänglich meist ein als eher diffus empfundener, heftiger Schmerzzustand im Bereich der
Nabelgegend\gACh{GABS}{2009-05-07}{Magengegend --> Nabelgegend}, der
innerhalb weniger Stunden in den rechten Unterbauch wandert. Übelkeit, Brechreiz und
\gACh{GABS}{2009-05-07}{Obstipation --> Verstopfung} können das Krankheitsbild
begleiten. Meist ist die Körpertemperatur bis 38\textcelsius  leicht erhöht
\gACh{GABS}{2009-05-07}{gelöscht da unerheblich für RS: , wobei eine rektale und
axilläre Messung eine Differenz um ca. 1\textcelsius aufweist}. 


\gACh{GABS}{2009-05-07}{alt: Unbehandelt kann der weitere Verlauf durch eine
Abszeßbildung und Perforation in
die Bauchhöhle gekennzeichnet sein. neu: Unbehandelt kann es zu einer Eiterung
und einem Durchbruch in die Bauchhöle kommen.} Unbehandelt kann es zu einer
Eiterung und einem Durchbruch in die Bauchhöle kommen. In diesem Stadium entwickelt sich ein
\gACh{GABS}{2009-05-07}{alt: Schockzustand, der eine absolute Lebensbedrohung
darstellt; neu: akutes Abdomen, siehe \ref{topic:akutes-abdomen}} akutes
Abdomen, siehe \ref{topic:akutes-abdomen}. 

Der Verlauf einer akuten Appendizitis ist
jeweils durch individuelle Faktoren wie Lebensalter und Allgemeinzustand des
Patienten geprägt. } {\gTabHeading{Symptome}

\begin{itemize}
\item Schmerzen im re. Unterbauch
\item \gE{Druck}schmerzen im re. Unterbauch
\item \gE{Loslaß}schmerz li. Bauch
\item Abwehrspannung 
\item \"Ubelkeit, Erbrechen
\item evt. Fieber
\end{itemize}}






\subsubsection{Gallenkolik}

\gDefinition{Kolik-artige Schmerzen infolge einer Blockade
der Gallenwege durch Gallensteine}

\gMarried 
{Gallensteine verursachen eine Verstopfung des
Gallengangs und foglich einen Rückstau von
Galleflüssigkeit in die Leber. Damit verbunden sind
starke, auf- und abschwellende, krampfartige
(kolikartige) Schmerzen im rechten Oberbauch,
klassischerweise mit Ausstrahlung in die rechte Schulterregion. Begleitendes Fieber ist möglich.

Der rechte Oberbauch ist äußerst druckschmerzhaft, der
Pateint ist sehr unruhig und agitiert
}{Abflu{\ss}stau der Gallenwege

R\"uckstau

Gallenblase kontrahiert gegen Widerstand (Stein)

\gTabHeading{Symptome}

\begin{itemize}
\item auf- und abschwellende, krampfartige
(kolikartige) Schmerzen im re. Oberbauch
\item Besonders nach fettreichen Mahlzeiten
\item evt. Gelbfärbung des Augenweiß (Gelbsucht (Ikterus))
\item evt. Übelkeit, Erbrechen
\end{itemize}}




\subsubsection{\"Osophagusvarizenblutung}

\gDefinition{Blutende "`Krampfadern"' in der
Speiseröhre.}\nopagebreak[4]

\gMarried
{Bei Patienten mit einer fortgeschrittenen Lebererkrankung
(Leberzirrhose, \ref{topic:leberzirrhose}) kann es zur
Bildung von Krampfader-artigen Gefäßen in der Speiseröhre
kommen. Diese Gefäße können massive Blutungen verursachen
die tödlich enden können. 

Je anch Stärke  äußern sich die Blutungen durch
schwallartiges Erbrechen von Kaffeesatz-artigen Mageninhalt (wenn das Blut von der
Magensäure bereits verändert wurde) oder von rotem,
frischen Blut. 

\gFazit{Der Patient ist aufgrund des Blutverlustes akut
vital bedroht, es besteht Schockgefahr.}}
{Oft bei Lebererkrankungen (Leberzirrhose)

Hoher Blutverlust möglich -- Lebensgefahr !

\gTabHeading{Symptome}
\begin{itemize}
\item Bluterbrechen im Schwall (Hämatemesis)
\end{itemize}}


\begin{gMassnahmenEnv}
    \paragraph{Spezielle Maßnhamen bei Ösophagusvarizenblutung}~
    
    \begin{itemize}
    \item \gTxMassVitBes
    \begin{itemize}
        \item $O_2$ 10-15\gLMin
        \item Lagerung: \gAuthNotes{Lagerung? Beine hoch?
        OK hoch?}
    \end{itemize}
    \item Allgemeine Maßnahmen der Schockbekämpfung
    ($\rightarrow$\,\ref{topic:schockbekaempfung})
    \item Spezielle Maßnahmen bei Hypovolämischen Schock
    ($\rightarrow$\,\ref{topic:schock-hypovolaemischer})
    \item \gZ{Notarzt: evt. Sengstaken-Blakemore-Sonde}
    \item Hospitalisation: Abteilung für Chirurgie, Voranmeldung
\end{itemize}
\end{gMassnahmenEnv}






\subsubsection{Magen-Darm-Grippe -- Gastroenteritis}

\emph{Syn.: Brechdurchfall, \gZ{Gastroduodenitis}}

Die Magen-Darm-Grippe ist eine chronisch oder akut
verlaufende Entzündung der Magen- und der Dünndarmschleimhaut. 

\gMarried
{Die Gastroenteritis wird durch Erreger
wie Bakterien oder Viren hervorgerufen. Es gibt völlig
harmlose aber auch lebensbedrohliche Verlaufsformen, dies ist abhängig vom Erreger und
letztendlich auch vom Alter des Patienten. Durch den
Durchfall und das Erbrechen kann es zu einem maßgeblichen Verlust von Wasser un Elektrolyten
kommen. Wenn der Patient
nicht ausreichend trinken kann besteht die Gefahr der Austrocknung (Exsikkose).

Besonders betroffen davon sind Kleinkinder, da sie
normalerweise einen höheren Wasseranteil am Körpergewicht
haben.  Die Symptome sind klassich: Oberbauchkrämpfe,
die ziemlich heftig sein können, Übelkeit, Erbrechen, Durchfälle.
}{~

\gTabHeading{Ursachen}: Vielfältig

\gTabHeading{Symptome}

\begin{itemize}
\item Erbrechen (Emesis)
\item Durchfall (Diarrhoe)
\item Wasser-/Elektrolytverlust
\end{itemize}

\gTabHeading{besonders gefährdet}

\begin{itemize}
\item (Klein-)Kinder
\item alte Menschen
\end{itemize}

Exsikkose-Gefahr!}






\subsubsection{Flüssigkeitsmangel -- Exsikkose}

\emph{Syn.: Dehydratation}

\gMarried
{\gAuthNotes{Fließtext fehlt}}
{Ursachen
\begin{itemize}
\item Fl\"ussigkeitszufuhr ${\downarrow}$
\item Fl\"ussigkeitsverlust ${\uparrow}$
\item forcierte Diurese  (Entwässerungsmedikamente)
\item starkes Schwitzen
\end{itemize}
Symptome
\begin{itemize}
\item RR ${\downarrow}$
\item stehende Hautfalten
\item trockene, belegte Zunge
\item Elektrolytentgleisung
\item Verwirrtheit, Schwindel, Apathie
\end{itemize}}



\subsection{Weniger häufige oder nicht-so-gut-erkennbare Erkrankungen}



\subsubsection{Bauchaortenaneurysma}

Zerrei{\ss}ung der Bauchaorta / Platzen  einer Aussackung
(Schwachstelle)

\gMarried
{\gAuthNotes{Fließtext fehlt!}}
{\gTabHeading{Symptome}

\begin{itemize}
    \item heftigste Schmerzen
("`Vernichtungsschmerz"')
\item     hoher Blutverlust ${\rightarrow}$ schnell im 
Schock!
\item evt. bereits bekanntes Aneurysma
${\rightarrow}$\,Anamnese
\item evt. tastbare pulsierende Schwellung
\end{itemize}
}

Oft ist bei den Patienten das Aneurysma bereits bekannt und
war noch nicht groß genug um behandelt zu werden
(${\rightarrow}$\,Anamnese!). Ansonsten ist es schwer das
Aneurysma präklinisch zu diagnostizieren. Im Vordergrund
steht der Volumenverlust und der entstehende hypovolämische
Schock. 



\subsubsection{Magen- und Zwölffingerdarmgeschwür}

\gZ{\emph{Syn. Magengeschwür: Gastritis}}

\gDefinition{Magengeschwür: Entz\"undung der
Magenschleimhaut mit Substanzdefekt.}

\gMarried
{\gAuthNotes{Fließtext fehlt: insbes. \underline{Wie} ist
der Schmerz, wann}

Es kommt zu Appetitlosigkeit, Völlegefühl, Übelkeit. Wenn
gleichzeitig eine Blutung besteht, dann sind Teerstühle
oder Kaffeesatz-artiges Erbrechen zu beobachten.}
{\gTabHeading{Ursachen}
\begin{itemize}
\item Autoimmunerkrankungen
\item bakteriell (Helicobacter pylori )
\item chemisch
\item Gallensäurenreflux
\item Medikamente (NSAR)
\end{itemize}

\gTabHeading{Symptome}

\begin{itemize}
\item Schmerzen (Oberbauch), \"Ubelkeit, Erbrechen
\item bei Blutung: Kaffesatz-Erbrechen,  Teerstuhl
\end{itemize}}





\subsubsection{Hepatitis}

\gDefinition{Entz\"undung der Leber}

Die infektiöse (virale) Hepatitis wird im Kapitel Hygiene
(\ref{topic:hepatitis-virale}weiter ausgeführt.

\gMarried
{\gAuthNotes{Fließtext fehlt}}
{\gTabHeading{Ursachen}

\begin{itemize}
    \item Gallen-Abflu{\ss}stau
    \item Vergiftungen Lebensmittel(Pilze)~/
    Medikamente(Paracetamol)~/ chron. Alkoholismus
    \item Virusinfektionen (Hep A, B, C, \ldots )
\end{itemize}

\gTabHeading{Symptome}

\begin{itemize}
    \item uncharakteristische Allgemeinsymptome (grippale Symptome,
Appetitlosigkeit)
\item Schmerzen im re. Oberbauch
\item Gelbfärbung Augen/Haut (Ikterus )
\end{itemize}}


Bei einem chronischen
Verlauf kommt es zu einem Umbau von
funktionierendem Lebergewebe zu
funktionslosem Bindegewebe. Man spricht
dann von der
\gT{Leberzirrhose}\label{topic-leberzirrhose}\index{Leberzirrhose}.




\subsubsection{Bauchfellentzündung}

\gZ{\emph{Syn.: Peritonitis}}

\gDefinition{Entz\"undung und Keimbesiedelung des 
Bauchfells}

\gMarried
{Durch die Infektion und Entzündung des Bauchfelles kommt es
zu Symptomen eines Akuten Abdomens. Ursächlich ist oft die Perforation eines Hohlorganes  (Magendurchbruch,
Darmperforation, Blinddarmdurchbruch, \ldots). 
} {\gTabHeading{Ursache}

\begin{itemize}
\item Perforation eines Hohlorganes
\item Infektion \"uber Blut- / Lymphwege oder eingewandert
über Fremdkörper
\end{itemize}

\gTabHeading{Symptome}

\begin{itemize}
\item Symptome eines Akuten Abdomens
\item Schmerzen (diffus)
\item evt. reflektorischer Darmverschluß
\item evt. Sepsis- und/oder Schockzeichen
\end{itemize}}



\gZ{
\subsection{Sonstige -- nicht akute -- Erkrankungen}
\subsubsection{Chronisch entzündliche Darmerkrankungen}
}



%%%%%%%%%%%%%%%%%%%%%%%%%%%%%%%%%%%%%%%%%%%%%%%%%%%%%%%%%%%
%%%
%%%%%%%%%%%%%%%%%%%%%%%%%%%%%%%%%%%%%%%%%%%%%%%%%%%%%%%%%%%

\subsubsection{Abschnitt oder Unterabschnitt}

\gDefinition{Defnitionstext --- Lorem ipsum dolor sit amet, consectetuer adipiscing elit, sed diam nonummy nibh euismod tincidunt ut laoreet dolore magna aliquam erat volutpat. Ut wisi enim ad minim veniam, quis nostrud exerci tation ullamcorper suscipit lobortis nisl ut aliquip ex ea commodo consequat.}

\gEs{Einleitungstext} --- Duis autem vel eum iriure dolor in
hendrerit in vulputate velit esse molestie consequat, vel
illum dolore eu feugiat nulla facilisis at vero eros et
accumsan et iusto odio dignissim qui blandit praesent
luptatum zzril delenit augue duis dolore te feugait nulla
facilisi. Nam liber tempor cum soluta nobis eleifend option
congue nihil imperdiet doming id quod mazim placerat facer
possim assum. Typi non habent claritatem insitam; est usus
legentis in iis qui facit eorum claritatem. Investigationes
demonstraverunt lectores legere me lius quod ii legunt
saepius. Claritas est etiam processus dynamicus, qui sequitur
mutationem consuetudium lectorum.

\gMarried{

\gEs{Blumige Beschreibung} --- Mirum est notare quam littera
gothica, quam nunc putamus parum claram, anteposuerit
litterarum formas humanitatis per seacula quarta decima et
quinta decima. Eodem modo typi, qui nunc nobis videntur parum
clari, fiant sollemnes in futurum. Lorem ipsum dolor sit
amet, consectetuer adipiscing elit, sed diam nonummy nibh
euismod tincidunt ut laoreet dolore magna aliquam erat
volutpat. Ut wisi enim ad minim veniam, quis nostrud exerci
tation ullamcorper suscipit lobortis nisl ut aliquip ex ea
commodo consequat.

Duis autem vel eum iriure dolor in
hendrerit in vulputate velit esse molestie consequat, vel
illum dolore eu feugiat nulla facilisis at vero eros et
accumsan et iusto odio dignissim qui blandit praesent
luptatum zzril delenit augue duis dolore te feugait nulla
facilisi.}
{
\gEs{Schlagworte} ---

\gTabHeading{Symptome}

\begin{itemize}
    \item Punkt 1
    \item Punkt 2
    \begin{itemize}
        \item a
        \item b
        \item a + b
    \end{itemize}
    \item Lorem
    \item ipsum
    \item Dolor
    \item Amet
\end{itemize}

\gTabHeading{Gefahren und Komplikationen}

\begin{itemize}
    \item Littera
    \item Littera gothica
    \item Litterarum
    \item Parum clar
\end{itemize}

}

\gEs{Fußtext} --- am liber tempor cum soluta nobis eleifend
option congue nihil imperdiet doming id quod mazim placerat facer possim assum.
Typi non habent claritatem insitam; est usus legentis in iis
qui facit eorum claritatem. Investigationes demonstraverunt
lectores legere me lius quod ii legunt saepius. Claritas est
etiam processus dynamicus, qui sequitur mutationem
consuetudium lectorum. Mirum est notare quam littera gothica,
quam nunc putamus parum claram, anteposuerit litterarum
formas humanitatis per seacula quarta decima et quinta
decima. Eodem modo typi, qui nunc nobis videntur parum clari,
fiant sollemnes in futurum.

\begin{gMassnahmenEnv}
    \paragraph{Spezielle Maßnahmen}
    
    \begin{itemize}
    \item \gTxMassVitBes (gTxMassVitBes)
    \begin{itemize}
        \item Lagerung: Oberkörper hoch
        \item $O_2$: 12\gLMin
        \item Ursachenforschung
    \end{itemize}
    \item Claritas
    \item Investigationes
    \item seacula
    \item futurum
\end{itemize}
    
\end{gMassnahmenEnv}

\section{Urologische Notfälle}


\subsection{Nierenkolik}

\gDefinition{Blockade des Harnleiters durch Nierensteine
mit massiven kolikartigen Schmerzen aufgrund des Harnstaus.}

\gMarried{\vspace{-1.2em}\paragraph{Ursache}
Die Nierenkolik ist eine äußerst schmerzhafte
akute Erkankung. Sie entsteht dadurch daß Nierensteine den
Harnleiter blockieren und verstopfe un dadurch Harn gestaut
wird und den Harnleiter dehnt. 
}{\gTabHeading{Ursache}
\begin{itemize}
\item Nierenstein(e)
\end{itemize}}

\gMarried{\vspace{-1.2em}\paragraph{Symptome}
Die Folge sind massive,
kolikartige Schmerzen, in der Flanke bzw. im Unterbauch der
jeweiligen Seite. Im Gegensatz zu anderen abdominellen
Erkrankungen ist der Patient extrem unruhig und agitiert.
Oft läuft er gestikulierend und wehklagend durch die
Gegend. Mitunter ist eine entsprechende Vorerkrankung
(\gE{Nierensteine}bereits bekannt und kann erfragt werden. 
}{\gTabHeading{Symptome}
\begin{itemize}
\item R\"ucken-/Flankenschmerzen
\item ev. Unterbauchschmerzen bei  Stein-/Sandabgang
\item {\quotedblbase}der Stein wandert, der Schmerz wandert, der Patient
wandert{\textquotedblleft}
\item Blut im Harn (Hämaturie)
\end{itemize}}

\paragraph{Versorgung} Im Vordergrund bei der Versorgung
steht der rasche zügige Transport an eine Abteilung für Urologie. Wenn die
Schmerzen sehr massiv sind muss eventuell eine
Schmerztherapie durch den Notarzt erfolgen.


\begin{gMassnahmenEnv}
    \paragraph{Spezielle Maßnahmen bei Nierenkolik}
    \begin{itemize}
        \item Lagerung nach Wunsch des Patienten
        \item schonender Transport
        \item Bei unerträglichen Schmerzen Notarzt zur
        Schmerztherapie beiziehen
        \item Hospitalisation: Abt. für Urologie
    \end{itemize}
\end{gMassnahmenEnv}





\subsection{Akutes Harnverhalten}

\gDefinition{Abflu{\ss}behinderung des Harns aus der
Harnblase, verscheidene Ursachen sind möglich.}



Das akute Harnverhalten ist besonders bei Männern im besten
oder oder fortgeschrittenem Alter häufig anzutreffen. 

\gMarried{\vspace{-1em}\paragraph{Ursachen} Hauptursache  dafür ist eine
altersbedingte Vergrößerung der Prostata, die die Harnröhre abschnürt. Davon abgesehen
können auch alte Narben, Medikamente oder eine Störung der
Schließmuskel der Harnblase\footnote{zum Beispiel im Rahmen eines
Querschnitt-Syndroms nach Wirbelsäulenverletzungen} zu
Harnverhalten führen. }
{\begin{itemize}
    \item Lähmung (Querschnittlähmung)
    \item Schlie{\ss}muskel-Krampf
    \item Verlegung der Harnröhre
    \item Medikamente
\end{itemize}}

Je nach Füllungszustand der Harnblase kann das
Harnverhalten sehr schmerzhaft sein.

Die Maßnhamen beschränken sich auf einen zügigen Transport
an eine Abteilung für Urologie. Eine Notarzt-Nachforderung
ist kaum sinnvoll, das NEF selten geeignietes Material zur
Harnableitung mitführen. 


\paragraph{Spezielle Maßnahmen bei akutem Harnverhalt}

\begin{itemize}
    \item Z\"ugiger Transport in eine Urologische Ambulanz
    \item \gZ{vorsichtig ablassen \"uber Harnkatheter oder
    Blasenpunktion (Spital)}
\end{itemize}



\subsection{Niereninsuffizienz und Nierenversagen}

\gAuthNotes{Abkz. CNI}

Meistens kommt es im Rahmen anderer Krankheiten
(Diabetes mellitus, Hypertonie, erhöhte Blutfette \ldots) zu einer massiven
Schädigung der kleinen Nierengefäße und in weiterer Folge zum Nierenversagen. Da
die zuvor erwähnten Krankheiten sehr häufig sind ist in weiterer Folge die
Niereninsuffizienz ebenfalls häufig anzutreffen. 

Als Therapie die steht der
Medizin anfänglich eine medikamentöse Behandlung, später jedoch nur die
Nierentransplantation oder die \gE{Blutwäsche}
(\gT{Dialyse}) zur Verfügung. Bei der klassichen Dialyse
wird der Patient jeden zweiten oder dritten Tag für
einige Stunden an eine Maschine angesclossen welche sein
Blut reinigen kann. Die "`Dialysefahrten"' sind eine der
Hauptaufgaben des Krankentransportdienstes.

\gEs{Der Ausfall der normalen Nierenfunktion hat
Auswirkungen} auf den Wasserhaushalt, den
Elektrolythaushalt und den Säurebasehaushalt. Diese Funktionen der Niere sind zum Überleben elementar (Vitalfunktionen zweiter
Ordnung). Dialysepatienten müssen daher sehr darauf achten nicht zuviel zu
trinken. Weiters können \gE{Elektrolytstörungen} sogar
\gE{lebensbedrohliche Notfälle} nach sich ziehen (bis hin
zum Herzstillstand).

Es ist daher wichtig auf den
Dialysepatienten -- auch im "`normalem"' Krankentransport --
immer ein wachsames Auge zu haben.







% \subsection{Hodentorsion}
% 
% Verdrehung eine Hodens im Hodensack  ={\textgreater}
% 
% Abschn\"urung von Samenstrang und  Blutgefä{\ss}en
% Oft nach Eintritt in Pubertät, ev. {\textgreater} 20 J.
% 
% \paragraph{Symptome:}
% 
% {\mdseries
% akuter heftiger Schmerz}
% 
% {\mdseries
% Schwellung}
% 
% \paragraph{Therapie:}
% 
% hochlagern
% 
% umgehende urologisch-chirurgische  Intervention (ad URO!!!)




\clearpage
\section{Schwangerschaft und geburtshilfliche Notfälle}
{\mdseries\itshape
Michael Withofner, Dr. med. univ.}

\subsection{Weiblicher Zyklus}

Dauer: 28 Tage

Follikelreifung (Hormon: \"Ostrogen )

Eisprung (Ovulation )

Gelbkörperphase (Hormon: Progesteron )

2 Möglichkeiten:

\begin{enumerate}
    \item Einnistung ={\textgreater} Schwangerschaft (Gestation )
    \item keine Einnistung ={\textgreater} Regelblutung  (Menstruation )
\end{enumerate}

Schwangerschaft (Gestation)

Eizelle:

8-10 Std. nach Eisprung befruchtungsfähig

benötigt 4-6 Tage f\"ur Wanderung Eierstock- Eileiter-Uterus

befruchtet ={\textgreater} Einnistung

Placenta bildet sich: 

Nährstoff- / O2 -Versorgung des Fetus

Hormon: ${\beta}$-HCG (in Blut und Urin der Mutter nachweisbar)

\paragraph{Schwangerschaft -- Verlauf}

\gT{Embryo} : Frucht bis zur 12. SSW

\gT{Fetus} : Frucht ab 12. SSW bis zur Geburt

erste Kindsbewegungen: 18.-20.SSW

Lungenreife: ca. ab 26. SSW 

Ab der \gEs{SSW ~24} ist das Kind theoretisch
\"uberlebensfähig (und ggfs. reanimationsplflichtig).

\subparagraph*{Normale Schwangerschaftsdauer} Die normale Schwangerschaft
dauert \gEs{38-42 Wochen}, das sind rund Kalendermonate oder 10
Lunarmonate\footnote{\gZ{lat. luna: Mond. Eine Mondphase hat 28 Tage.}} Der
Fortschritt der Schwangerschaft wird in \gE{Schwangerschaftswochen}, Abk. SSW,
angegeben.

Die Schwangerschaft wird in 3 Drittel unterteilt, welche
\gT{Trimenon}\index{Trimenon} heißen.

Von einer \gEs{Fr\"uhgeburt} spricht man wenn die Geburt bis zur vollendeten 37.
SSW erfolgt.





\subsection{Geburt}

\paragraph{Verlauf}~

Die Geburt verläuft in \gE{3} Phasen:

%\begin{gWideParEnv}
\begin{center}
\begin{tabulary}{\linewidth}{LLL}
\toprule
Phase
 &
Beschreibung
 &
Wehen
\\\midrule
\gT{Eröffnungs\-phase}
 &
Eröffnungswehen

Gebärmutterhals (Cervix) verk\"urzt sich,  Muttermund
öffnet sich

Blasensprung

\gE{Liegender Transport}, sonst Beschleunigung der  Geburt
und Gefahr eines  \gEs{Nabelschnurvorfalls}!! 

Kindskopf tritt in Geburtskanal ein
 &
10-15 min.
\\\hline
\gT{Austreibungs\-phase}
 &
Pre{\ss}wehen

Pausen wichtig (Erholungsphase f\"ur Fetus)

Muttermund vollständig eröffnet

\gT{Pre{\ss}wehen}: Kontraktion Bauch- plus 
Uterusmuskulatur

Kopf wird sichtbar -- normalerweise  Hinterhauptslage (alle
anderen Lagen stellen  Komplikationen dar)

\gT{Dammschutz} ${\rightarrow}$ sonst
Dammri{\ss}, Schutz des Kindes vor Stuhl der Mutter

Kopf geboren ${\rightarrow}$ Körper folgt schnell nach
${\rightarrow}$  auffangen!

Absaugen ${\rightarrow}$ abnabeln ${\rightarrow}$ warm
einpacken ${\rightarrow}$ der Mutter auf die Brust legen
&
2-3 min.
\\\hline
\gT{Nachgeburts\-phase}
 &
Nach 10-20 min. neuerliche Wehen

Geburt der Placenta  (Nachgeburt)
 &
\\\bottomrule
\end{tabulary}
\end{center}
%\end{gWideParEnv}

\subparagraph{Neugeborenes - Vitalitätsbeurteilung}
\subparagraph*{APGAR-Score}

\begin{enumerate}
    \item A ussehen
    \item P uls
    \item G esichtsbewegung
    \item A ktivität
    \item R espiration
\end{enumerate}

je 0-2 Pkt., jeweils nach 1min/5min/10min

\subsection{Notfälle - Fr\"uhschwangerschaft}

\subsubsection{Eileiterschwangerschaft -- Tubaria}

Fehlgeburt (Abortus )

Notfälle - Tubaria

= extrauterine Schwangerschaft

Embryo hat sich in Eileiter (Tuba ) eingenistet

ab gewisser Grö{\ss}e reicht Platz nicht mehr aus

Tube rei{\ss}t / platzt


Blutung, Schmerz, akutes Abdomen, Schock

{\mdseries
${\rightarrow}$ Bei Frauen mit akutem Abdomen immer nach möglicher
Schwangerschaft fragen!}

\subsubsection{Abortus}
{\mdseries
Beendigung einer Schwangerschaft vor Beendigung der 28.SSW}

{\mdseries
Fr\"uhgeburt: 24.-37. SSW}

{\mdseries
Feten {\textgreater}24.SSW haben heute  \"Uberlebenschance!! (dank 
neonatologischer Intensivmedizin)}

\paragraph{Symptome}

\begin{itemize}
\item vaginale Blutung
\item Abgang von Klumpen
\item Wehenartige Schmerzen
\end{itemize}

\paragraph{Therapie}

\begin{itemize}
\item sterile Vorlagen
\item Lagerung nach Fritsch
\item falls erforderlich Schocktherapie
\item je nach Zustand
\item NAW
\item direkt Gyn
\end{itemize}


\subsection{Notfälle -- Spätschwangerschaft}

\subsubsection{Vorzeitige Plazentalösung}
{\mdseries
Aufgrund einer Einblutung zwischen Plazenta und Gebärmutter kann es zu
einer vorzeitigen Ablösung der Plazenta mit dann ungehinderter
Blutung kommen.}

\begin{itemize}
\item {\mdseries
\gEs{Lebensgefahr} f\"ur Mutter und Kind! }

\end{itemize}
\begin{itemize}
\item Blutung der Mutter und Mangelversorgung des Kindes 
\end{itemize}

Die Diagnose ist durch en Sanitäter nicht sicher zu stellen. Im
Vordergrund steht daher die Behandlung der Symptome (Schock) und ein
z\"ugiger Transport. 

\paragraph{Symptome}

\begin{itemize}
\item unspezifisch:

\begin{itemize}
\item Schock
\item akutes Abdomen
\end{itemize}
\end{itemize}

\begin{gMassnahmenEnv}
    \paragraph{Spezielle Massnahmen}
    \begin{itemize}
        \item \gTxMassVitBes
        \item Schockbekämpfung
        \item keine vaginale Untersuchung /  Scheidentamponade
        \item Z\"ugiger Transport mit Voranmeldung
    \end{itemize}
\end{gMassnahmenEnv}


\subsubsection{Vena-cava-Syndrom}

Es kann vorkommen dass das \gEs{Kind die untere
Hohlvene der Mutter abdr\"uckt}. Dadurch wird der Blutstrom zum Herzen
vermindert und es kann zu einem \gEs{plötzlichen
Kollaps der Mutter }und einer Unterversorung des Kindes kommen. Die
Therapie besteht in einer Linksseitenlagerung, da die Hohlvene eher
rechts der Körpermitte verläuft. 

{\mdseries
\gEs{Jede hochschwangere Pateintin }sollte daher in
Linksseitenlage und nicht in R\"uckenlage transportiert werden. }

\paragraph{Ma{\ss}nahmen}

\begin{itemize}
    \item Linksseitenlage (Hohlvene verläuft eher rechts)
\end{itemize}

\subsubsection{EPH-Gestose -- Präeklampsie - Eklampsie}
\label{topic:eph-gestose}\index{EPH-Gestose}\index{Präeklampsie}\index{Eklampsie}\index{Krampfanfall!eklamptischer}

Synonym: \gT{Präeklampsie}

Die \gT{EPH-Gestose}, auch \gT{Präeklampsie} genannt,
ist eine der gefüchtetsten Komplikationen der Spätschwangerschaft.
Unbehandelt kann sie zum einem \gE{eklamptischen
Krampfanfall} (\gT{Eklampsie}, ähnlich einem "`normalen"'
zerebralen Krampfanfall) unf in weiterer Folge zum Tod von Mutter und Kind führen. 

\gFazit{Die Präeklampsie (EPH-Gestose) kann zu einer
Eklampsie (lebensgefährlicher Krampfanfall) führen.}

\paragraph{Symptome}

\begin{enumerate}
    \item Edema (\"Odeme)
    \item Proteinuria (Proteinurie = Eiwei{\ss} im Harn)
    \item Hypertonus (RR  ${\uparrow}$)
\end{enumerate}


RR {\textgreater}140/90 bei einer Schwangeren ist zu  hoch!!!

Mutter-Kind-Pa{\ss} Untersuchungen zur  Fr\"uherkennung



\begin{gMassnahmenEnv}
  \paragraph{Spezielle Maßnahmen bei Verdacht auf
Präeklampsie}
\begin{itemize}
\item schonender Transport
\item \"Uberwachung, O2
\item Notarzt: Krampfdurchbrechung, RR-Senkung
\item (im Vollbild: vorzeitige Beendigung der  Schwangerschaft)
\end{itemize}  
\end{gMassnahmenEnv}

\begin{gMassnahmenEnv}
  \paragraph{Spezielle Maßnahmen bei Eklampsie --
  Eklamptischer Krampfanfall}
Grundsätzlich ist zu verfahren wie bei jedem anderen
Krampfanfall, vgl. Kap. \ref{topic:krampfanfall}. Bedenke:
Der eklamptische Krampfanfall ist für Mutter und Kind lebensgefährlich!
\end{gMassnahmenEnv}


\subsubsection{Vorzeitiger Fruchtwasserabgang}

Der Fetus schwimmt während  Schwangerschaft in steriler
Fl\"ussigkeit (Fruchtwasser). Bei vorzeitige Eröffnung der
Fruchtblase ergeben sich für ihn die Gefahr
eines \gEs{Nabelschnurvorfalls} oder einer \gEs{Infektion}
(Keimbesiedlung der Fruchthöhle). Die Patientin ist an
einer geburtshilflichen Abteilung vorzustellen.

\begin{gMassnahmenEnv}
    \paragraph{Ma{\ss}nahmen}

\begin{itemize}
\item Lagerung: Linksseitenlagerung, Becken erhöht
(Schwerkraft verzögert Geburt), Lagerung nach Fritsch
\item Nicht mehr aufstehen lassen! Liegender Transport!
\item Hospitalisation: Abteilung für Gynäkologie und
Geburtshilfe, Kreißsaal
\end{itemize}
\end{gMassnahmenEnv}



\subsection{Notfälle während der Geburt}

\subsubsection{Nabelschnurvorfall}

\paragraph{Gefahren:}

Abklemmung der Nabelschnur  (Sauerstoffversorgung  [F0EA?] )

Strangulation (Gehirnperfusion [F0EA?] )

\paragraph{Therapie:}

Linksseitenlagerung, Beckenhochlagerung

Mutter soll NICHT pressen (Geburt w.m. erst  im KH!!!)

Wehenhemmung (hecheln, medikamentös) 

Nabelschnurumschlingung des  kindlichen Halses während der Geburt 
sofort lösen!

Placenta praevia

Plazenta liegt unterhalb des Fetus in  Geburtskanal

Geburtshindernis

massive Blutungen

\paragraph{Ma{\ss}nahmen}

Lagerung nach Fritsch

Schockbekämpfung

NAW

im KH: Notsectio

\subsubsection{Pathologische Geburtslagen}

alle Lagen, bei denen nicht zuerst der  Kopf im Geburtskanal erscheint 

Stei{\ss}lage

Beckenendlage

Armvorfall

\paragraph{Gefahr}

Abdr\"ucken der Nabelschnur

Geburtsstopp

Geburt nur im KH durchf\"uhrbar (Sectio- Bereitschaft!)

In jedem Fall Geburt  verhindern/hinauszögern!

\subsubsection{Uterusatonie}

normalerweise zieht sich in der  Nachgeburtsphase die Uterusmuskulatur 
zusammen

hier: Erschlaffung der Uterusmuskulatur  nach Plazentalösung
={\textgreater} gro{\ss}flächige  Blutung!

\paragraph{Ma{\ss}nahmen}

\begin{itemize}
    \item Allgemeine Ma{\ss}nahmen bei lebensbedrohlichen Zuständen
    \item Schockbekämpfung
    \item Lagerung: Fritsch, Beine hoch
    \item Evt. Druck mit beiden Fäusten oberhalb  Schambein ${\rightarrow}$
    Blutungsverminderung
    \item Z\"ugiger Transport mit Voranmeldung
\end{itemize}


\subsubsection{Uterusruptur}
{\mdseries
plötzlicher Schmerz während Wehen,  dann keine Wehen mehr
${\rightarrow}$ Geburtsstop}

\gFazit{Lebensgefahr f\"ur Mutter (Blutung) und  Kind (O2
-Mangel!)}

\paragraph{Therapie}

\begin{itemize}
    \item Allgemeine Ma{\ss}nahmen bei lebensbedrohlichen Zuständen
    \item Schockbekämpfung
    \item Z\"ugiger Transport mit Voranmeldung
\end{itemize}

\paragraph{Asphyxie}

Neugeborenes atmet nicht oder nicht  ausreichend

siehe APGAR - Score

\paragraph{Therapie}

absaugen

Fu{\ss}sohlen und R\"ucken klopfen, reiben

O2

Reanimation, Notarzt


\subsection{Sonstige gynäkologische Erkrankungen und
Notfälle}

\subsubsection{Vaginale Blutungen}

\gMarried
{Die vaginale Blutung ist "`im Feld"' schwer abzuklären und
kann viele Ursachen haben. Wichtig ist festzuhalten daß
jede vaginale Blutung abgeklärt gehört, da sich dahinter
auch ein Tumorgeschehen (besonders bei Patientinnen in der
Menopause) verstecken kann. Darüber hinaus kommt eine
verstärkte Monatsblutung, Verletzung (nach Sturz, 
Geschlechtsverkehr oder nach Vergewaltigungen),
Defloration\footnote{Entjungferung} oder Blutungen im Rahmen
einer Schwangerschaft in Betracht. Bei vaginalen Blutungen in Kindern darf man \gE{nicht}
automatisch von einem erfolgten Mißbrauch ausgehen. }
{\begin{itemize}
    \item verstärkte Monatsblutung
    \item Verletzungen (Sturz, Geschlechtsverkehr, Vergewaltigung,
\ldots)
\item Defloration
\item Tumor
\item Im Rahmen einer Schwangerschaft
\end{itemize}}

Eine akute vitale Bedrohung besteht seltenst, dennoch muß
sie routinemäßig vor allem anhand des Blutverlustes und der
Vitalwerte eingeschätzt werden. 


\begin{gMassnahmenEnv}
    \paragraph{Spezielle Maßnahmen bei vaginaler Blutung}

\begin{itemize}
\item Lagerung nach Fritsch
\item Hospitalisation: Abteilung für
Gynäkologie, bei Wehen Kreißsaaal
\end{itemize}
\end{gMassnahmenEnv}





\clearpage
%%%%%%%%%%%%%%%%%%%%%%%%%%%%%%%%%%%%%%%%%%%%%%%%%%%%%%%%%%%
%%% Pädiatrie                                           %%%
%%%%%%%%%%%%%%%%%%%%%%%%%%%%%%%%%%%%%%%%%%%%%%%%%%%%%%%%%%%

\section{Pädiatrie - Kinderheilkunde}

\author{Michael Withofner, Dr. med. univ.}

\subparagraph*{Prinzipieller Umgang mit Kindern}

\begin{itemize}
\item schonend
\item jegliche Aufregung vermeiden
\item Bezugsperson wichtig
\item Beruhigung
\item Kind während Transport am Besten eng bei  Vertrauensperson belassen
\end{itemize}





\subsection{Anatomische und physiologische  Besonderheiten}

Kinder sind keine kleinen  Erwachsenen!

\begin{center}
\begin{tabulary}{12cm}{LL}\toprule
a
&
b
\\\midrule
Atemfrequenz:
&
{\mdseries Neugeborenes: 40x / min.}

{\mdseries Säugling: 20x / min.}

{\mdseries Kleinkind: 15x / min.}

{\mdseries Beatmung beim Kind immer in  Neutralposition (Kopf nicht
\"uberstreckt)}

{\mdseries Kehlkopf steht höher, ist verkippt}

\\\hline
Zeichen der Atemnot
&
Einziehungen zwischen Rippen

\index{Nasenfl\"ugeln}\gEs{Nasenfl\"ugeln}
beschleunigte Atmung
Atemgeräusche
\\\hline
Pulstasten beim Säugling:
&
Oberarmarterie (Oberarm Innenseite), Achsel
\\\hline
Herzfrequenz:
&
{\mdseries Neugeborenes: 140 -- 160 / min.}

{\mdseries Säugling: 120 / min.}

{\mdseries Kinder: 80-100 / min.}

{\mdseries Bradykardie  ist fast immer  hypoxiebedingt ={\textgreater}
lebensgefährlich!!!}

\\\hline
Wasser-, Elektrolyt- und Wärmehaushalt
&
{\mdseries Kinder haben im Verhältnis zu ihrem  Körpergewicht eine
grö{\ss}ere Körperoberfläche  als Erwachsene!}

{\mdseries höherer Wasseranteil}

{\mdseries Austrocknung und Unterk\"uhlung schneller  möglich!!!}
\\\bottomrule
\end{tabulary}
\end{center}


\subsection{Fremdkörperaspiration }

\gDefinition{Verlegung der Atemwege durch Fremdkörper wie
zum Beispiel Nahrungsmittel, Spielzeug o.ä.}

\paragraph{Symptome}

\begin{itemize}
    \item plötzlich starker Husten
    \item Atemnot
    \item Stridor oder Keuchen
    \item W\"urgen
\end{itemize}

\begin{gMassnahmenEnv}
    
\subparagraph*{Spezielle Maßnahmen}

\begin{itemize}
    \item \gTxMassVitBes
    \begin{itemize}
        \item $O_2 15\,l$
        \item Lagerung: aufrecht, Oberkörper hoch,
        Atemhilfsmuskulatur unterstützen
    \end{itemize}
    \item Schlag zwischen die Schulterblätter tangential  (Kopf tief)
    \item Bei grö{\ss}eren Kindern auch Heimlich-Maneuvre
    \item keine Entfernungsversuche
\gAuthNotes{Mit ERC-Guidelines abstimmen !!}
    \item Beatmung bei Atemstillstand
\end{itemize}

\end{gMassnahmenEnv}

\subsection{Laryngitis}

Pseudokrupp / Kruppsyndrom

Behinderung der oberen Atemwege

durch Infektion mit verschiedenen Viren

{\mdseries
Glottis und Subglottis}

\paragraph{Symptome:}

\begin{itemize}
    \item Pfeifendes Geräusch beim Einatmen
    \gZ{[inspiratorischer Stridor)}
    \item trockener, bellender Husten
    \item Heiserkeit
    \item Dyspnoe
    \item kaum Fieber
    \item typisches Alter: 3 Monate -- 6 Jahre
\end{itemize}

\paragraph{Therapie:}

\begin{itemize}
    \item kalte  Luft oder hei{\ss}er  Dampf
    \item kaltes Wasser trinken lassen (DD Epiglottitis)
    \item Notarzt: Cortison o.ä.
    \item immer Eltern und  Kind beruhigen
\end{itemize}



\subsection{Epiglottitis}

bakterielle Entz\"undung und daraus  folgende Schwellung des
Kehldeckels

selten, aber eventuell lebensbedrohlich

4.-5.Lj. Gehäuft

schwere Erkrankung, mit richtiger  Behandlung aber meist gute Prognose

\paragraph{Symptome:}

\begin{itemize}
\item {\mdseries
hochfieberhafte Erkrankung}

\item {\mdseries
{\quotedblbase}hot potato voice{\textquotedblleft}}

\item {\mdseries
inspiratorischer Stridor}

\item {\mdseries
Mund offen, Speichelflu{\ss}, spricht kaum!!!}

\item {\mdseries
Angst!!!!!}

\end{itemize}
\subparagraph*{
Therapie:}

\begin{itemize}
\item {\mdseries\color[rgb]{0.0,0.0,0.5019608}
Notarzt}

\item {\mdseries\color[rgb]{0.0,0.0,0.5019608}
O2 , Luftbefeuchtung}

\item {\mdseries\color[rgb]{0.0,0.0,0.5019608}
Kinder sitzen lassen!}

\end{itemize}


\subsection{Laryngitis vs. Epiglottitis}

{\mdseries
Laryngitis:}

{\mdseries
wenig bedrohlich}

{\mdseries
Kind ist weinerlich}

{\mdseries
kann schlucken}

{\mdseries
schreit ev.}

\subsection{SIDS -- Sudden Infant Death Syndrome -- Plötzlicher Kindstod}

{\mdseries
{\quotedblbase}\index{plötzlicher
Kindstod}\gT{plötzlicher Kindstod}{\textquotedblleft}}

\begin{itemize}
\item plötzlich, meist im Schlaf

\item vorwiegend 1.-4.Lebensmonat

\item Saisonal (Jän bis März)

\item familiär und bei Fr\"uhchen gehäuft

\end{itemize}
\subparagraph*{
Ursachen}

\begin{itemize}
\item letztendlich unbekannt

\item zentrale Atemregulationsstörung (unreifes  Atemzentrum?)

\end{itemize}
{\mdseries
Warnzeichen:  lange Atempausen}

{\mdseries
Prophylaxe: Apnoemonitoring bei  Risikokindern}

\paragraph{Vorgehensweise}

{\mdseries
Eltern sollen handeln können: Gef\"uhl geben,  da{\ss} sie alles f\"ur
ihr Kind getan haben (sonst:  Schuldgef\"uhle)}

{\mdseries
im Zweifelsfall Pseudoma{\ss}nahmen (Transport  ins KH unter Reanimation
statt Belassung)}

{\mdseries
Zwilling unbedingt hospitalisieren!}

\subsection{Ertrinkungsunfall}\index{Ertrinkungsunfall}

{\mdseries
Tod in den ersten 24 h nach Unfall: Ertrinken}

{\mdseries
ansonsten:
\index{Beinahe-Ertrinken}\gT{Beinahe-Ertrinken}}

{\mdseries
Kleinkinder sind gefährdet (mangelnde  Aufsicht, können nicht
schwimmen)!!!!!}

\subparagraph*{
Faktoren:}

\begin{itemize}
\item {\mdseries
Kleinkind (kann noch nicht schwimmen,  Panik!)}

\item {\mdseries
Gefahrenquellen:}

\item {\mdseries
Gartenteiche, Swimmingpools, schlecht  gesichert}

\item {\mdseries
Gewässer}

\item {\mdseries
unbeaufsichtigte Wannenbäder}

\item {\mdseries
mangelnde Aufsicht}

\end{itemize}
\subparagraph*{
Therapie:}

\begin{itemize}
\item {\mdseries\color[rgb]{0.0,0.0,0.5019608}
Vor Ort: CPR}

\end{itemize}
\subparagraph*{
Nobody is dead until warm and dead!}

\subsection{Störungen des Nervensystems}

\subsubsection{Krampfanfälle:}
\begin{itemize}
\item {\mdseries
fokal}

\item {\mdseries
generalisiert}

\item {\mdseries
Gelegenheitskrämpfe:}

\item {\mdseries
Einfache Fieberkrämpfe}

\item {\mdseries
Komplizierte Fieberkrämpfe}

\item {\mdseries
Andere (Meningitis, Trauma, Synkope etc.)}

\item {\mdseries
Affektkrämpfe:}

\item {\mdseries
Pavor nocturnus}

\end{itemize}
\paragraph{Fieberkrämpfe}
\subparagraph*{
Einfach:}

\begin{itemize}
\item erbliche Vorbelastung

\item schlaff oder generalisiert krampfend

\item selbstlimitierend

\item reifungsbedingt (Kleinkinder {\textless} 5 J.), kein 
Dauerschaden, kein  \"Ubergang in eine  Epilepsie

\end{itemize}
\subparagraph*{
Kompliziert:}

\begin{itemize}
\item dauern länger

\item Bei Fieber öfters hintereinander

\item Bei primär neurologisch auffälligen Kindern  gehäuft

\item Schlechtere Prognose

\item Neigung, in Epilepsie \"uberzugehen

\item IMMER medikamentös unterbrechen  (Notarzt)

\item Krämpfe {\textgreater} 20 min. ={\textgreater} Dauerschaden!!! 

\end{itemize}
\subparagraph*{
Therapie:}

\begin{itemize}
\item {\mdseries\color[rgb]{0.0,0.0,0.5019608}
Wie {\quotedblbase}normaler{\textquotedblleft} Krampfanfall:}

\item \item {\mdseries\color[rgb]{0.0,0.0,0.5019608}
\gZ{(Fiebersenkung unter  38,5{\textdegree}C
(medikamentös,  mechanisch))}}

\item {\mdseries\color[rgb]{0.0,0.0,0.5019608}
\gZ{(Krampfdurchbrechung, falls notwendig 
medikamentös)}}

\end{itemize}


\subsection{Kindesmi{\ss}handlung}

{\mdseries
Absichtliche psychische oder physische  Gewalt gegen Kinder, die zu 
Verletzungen, Tod und  Entwicklungsstörungen f\"uhren kann}

{\mdseries
Kommt in ALLEN sozialen Schichten vor}

{\mdseries
Hinweise:}

{\mdseries
Langes Intervall zwischen  Verletzungszeitpunkt und Verständigung der 
Rettung}

{\mdseries
Häufiger Arzt-/KH-Wechsel}

{\mdseries
Anamnese und Verletzungsmuster  widersprechen einander!}

{\mdseries
Kindesmi{\ss}handlung}

\subparagraph*{
Symptome:}

{\mdseries
mehrere Verletzungen verschiedenen Alters  (Röntgen)}

{\mdseries
erkennbarer {\quotedblbase}verletzungsformender{\textquotedblleft} 
Gegenstand (Schuhabdruck etc.)}

{\mdseries
Sch\"uttelverletzungen (Subduralhämatom)}

{\mdseries
Zeichen der Vernachlässigung}

{\mdseries
Genitalregion (Verletzungen, Schmerzen ={\textgreater}  sexueller
Mi{\ss}brauch)}

{\mdseries
Verhaltensstörungen}

{\mdseries
Kindesmi{\ss}brauch}

\subparagraph*{Ma{\ss}nahmen:}

{\mdseries
Daran denken!!! (Verbrennungen, SHT,  Knochenbr\"uche, ...)}

{\mdseries
Vorsicht mit voreiligen  Verdachtsäu{\ss}erungen!!!}

{\mdseries
ev. Vorwand f\"ur Hospitalisierung suchen}

{\mdseries
dem KH Verdacht vorsichtig mitteilen}

\section{Psychiatrische Notfälle}

Michael Withofner, Dr. med. univ.

\dictum[Otto Pötzl, nach \cite{Jolesch2}]{Der, der den
Schlüssel hat, ist der Arzt.}
\cite{Jolesch2}
{\mdseries
10\% aller psychiatrischen Erkrankungen  liegt eine organische Ursache
zugrunde.}

{\mdseries
Bei Erstmanifestation einer  psychiatrischen Erkrankung MUSS  daher eine
organische Ursache  ausgeschlossen werden!!!}

{\mdseries
Der Umgang mit psychiatrischen Patienten ist heikel und erfordert viel
Feingef\"uhl sowie die Fähigkeit, auf das Gegen\"uber einzugehen. }

\subparagraph*{Akute Intervention}

{\mdseries
Indikation:}

{\mdseries
Unmittelbare Lebensgefahr (Eigengefährdung)}

{\mdseries
Gesundheit und Leben anderer gefährdet (Fremdgefährdung)}

{\mdseries
erschwerend: unkooperativer Patient}

\subparagraph*{2 Möglichkeiten:}

\begin{enumerate}
\item Freiwillige Behandlung (Einwilligung des Patienten)

\item Unterbringung (Behandlung gegen oder ohne Willen des Patienten)

\end{enumerate}
{\mdseries\color[rgb]{0.0,0.0,0.5019608}
Unterbringung (gem. Unterbringungsgesetz) }

{\mdseries
fr\"uher: {\quotedblbase}Parere{\textquotedblleft}}

{\mdseries
Die \index{Unterbringung}Unterbringung wird durch das
\index{Unterbringungsgesetz}Unterbringungsgesetz geregelt. Auszug:}

{\mdseries
{\S} 3. In einer Anstalt darf nur untergebracht werden, wer}

{\mdseries
1. an einer \textbf{psychischen Krankheit} leidet und im Zusammenhang
damit \textbf{sein Leben} oder seine \textbf{Gesundheit} oder das Leben
\textbf{oder} die Gesundheit \textbf{anderer} ernstlich und erheblich
gefährdet und}

{\mdseries
2. \textbf{nicht in anderer Weise}, insbesondere au{\ss}erhalb einer
Anstalt, ausreichend ärztlich behandelt oder betreut werden kann.}

{\mdseries
Unterbringung ohne Verlangen}

{\mdseries
{\S} 8. Eine Person darf gegen oder ohne ihren Willen nur dann in eine
Anstalt gebracht werden, wenn ein im \textbf{öffentlichen
Sanitätsdienst stehender Arzt} oder ein \textbf{Polizeiarzt} sie
untersucht und bescheinigt, da{\ss} die Voraussetzungen der
Unterbringung vorliegen. In der Bescheinigung sind im einzelnen die
Gr\"unde anzuf\"uhren, aus denen der Arzt die Voraussetzungen der
Unterbringung f\"ur gegeben erachtet.}

{\mdseries
{\S} 9. (1) Die \textbf{Organe des öffentlichen Sicherheitsdienstes}
sind berechtigt und verpflichtet, eine Person, bei der sie aus
besonderen Gr\"unden die Voraussetzungen der Unterbringung f\"ur
gegeben erachten, zur Untersuchung zum Arzt ({\S} 8) zu bringen oder
diesen beizuziehen. Bescheinigt der Arzt das Vorliegen der
Voraussetzungen der Unterbringung, so haben die Organe des
öffentlichen Sicherheitsdienstes die betroffene Person in eine
Anstalt zu bringen oder dies zu veranlassen. Wird eine solche
Bescheinigung nicht ausgestellt, so darf die betroffene Person nicht
länger angehalten werden.}

{\mdseries
(2) \textbf{Bei Gefahr im Verzug} können die \textbf{Organe des
öffentlichen Sicherheitsdienstes} die betroffene Person \textbf{auch
ohne Untersuchung und Bescheinigung in eine Anstalt bringen}.}

{\mdseries
(3) Der Arzt und die Organe des öffentlichen Sicherheitsdienstes haben
unter möglichster Schonung der betroffenen Person vorzugehen und die
notwendigen Vorkehrungen zur Abwehr von Gefahren zu treffen. Sie haben,
soweit das möglich ist, mit psychiatrischen Einrichtungen
au{\ss}erhalb einer Anstalt zusammenzuarbeiten und
\textbf{erforderlichenfalls den }\textbf{örtlichen Rettungsdienst
beizuziehen}.}

{\mdseries
Daraus ergibt sich bei
\gEs{fehlender}\gEs{
Freiwilligkeit}:}

{\color{black}
Voraussetzungen ({\S}3)}

\begin{itemize}
\item \gEs{Psychiatrische Erkrankung}
(ursächlich!), (${\rightarrow}$~{\S}3~Abs.~1)

\item \gEs{Eigen-}  und/oder
\gEs{Fremdgefährdung},
(${\rightarrow}$~{\S}3~Abs.~1)

\item \gEs{Keine weniger einschränkende
Möglichkeit} zur Abwendung der Gefährdung,
(${\rightarrow}$~{\S}3~Abs.~2)

\end{itemize}
{\color{black}
Vorgehensweise ({\S}{\S}~8, 9)}

Ein {\quotedblbase}im öffentlichen Sanitätsdienst stehender Arzt
oder ein Polizeiarzt{\textquotedblleft} die Voraussetzungen zu
bescheinigen und zu Begr\"unden. In Wien ist das der
{\quotedblbase}\index{Amtsarzt}\gEs{Amtsarzt}{\textquotedblleft},
welcher \"uber die Polizei verständigt wird. 

\gEs{Der Notarzt, auch wenn er von der Gemeinde Wien
gestellt wird, hat nicht die Kompetenz eine Unterbringung anzuordnen}.
(${\rightarrow}$~{\S}8)

F\"ur den Transport ist grundsätzlich die Polizei zuständig
(${\rightarrow}$~{\S}9~Abs.~1)! Die Polizei darf des Rettungsdienst
quasi als {\quotedblbase}Taxi{\textquotedblleft} beiziehen
(${\rightarrow}$~{\S}9~Abs.~3). 

Bei Gefahr im Verzug darf die Polizei auch ohne Bescheinigung eine
Person auf eine psychiatrische Abteilung zwecks Unterbringung
verbringen (${\rightarrow}$~Abs.~2). 

Im Zweifelsfall Polizei nachalarmieren...

{\color{black}
Weiters: }

\gEs{Eine Unterbringung kann nur in einer
}\gEs{psychiatrischen}\gEs{
Abteilung erfolgen!}

Eine Unterbringung \gEs{kann auch freiwillig}
erfolgen, da dann auch der Transport freiwillig stattfindet sind dabei
keine besonderen Bestimmungen zu beachten.

{\mdseries\color[rgb]{0.0,0.0,0.5019608}
Umgang mit Patienten}

{\mdseries
Eigenschutz!}

{\mdseries
ruhiger Kontakt}

\subparagraph*{
Situation:}

{\mdseries
Belassung? (ist Patient \"uberhaupt  reversfähig?)}

{\mdseries
Unterbringung? (notwendig / rechtlich  zulässig?)}

{\mdseries
Polizei / Amtsarzt bzw. Notarzt? (stabil?  Eskalation?)}

\subparagraph*{Leitsymptome}

\begin{itemize}
\item Verwirrtheit, Bewu{\ss}tseinseintr\"ubung
\item Realitätsverlust (Wahn, Halluzinationen )
\item psychomotorische Unruhe, Aggressivität
\item Depression , Verzweiflung, Suizidalität
\item Angst
\item Antriebslosigkeit, Stupor
\item Verwirrtheit / Bewu{\ss}tseinstr\"ubung
\item einfache Verwirrung
\item Altersdemenz
\item delirante Verwirrung
\item Einfache Verwirrung

\end{itemize}
{\mdseries
Kann bei allen {\quotedblbase}normalen{\textquotedblleft},
{\quotedblbase}körperlichen{\textquotedblleft} Notfällen, die mit
Bewusstseins\-ein\-schränkungen ein\-her\-gehen, auftreten!}

{\mdseries
Bsp.: Epilepsie , Commotio cerebri, Exsikkose etc.}

{\mdseries\color[rgb]{0.0,0.0,0.5019608}
Altersdemenz}

{\mdseries
Fortschreitende Einschränkung der  geistigen Leistungsfähigkeit des
Gehirns  im Alter}

\subparagraph*{
Betrifft}

\begin{itemize}
\item Gedächtnis
\item Denkvermögen
\item Sprache
\item Motorik
\item Persönlichkeitsstruktur
\item delirante Verwirrung
\end{itemize}

\paragraph{Symptome}


Desorientierung

zeitlich

örtlich

psychomotorische Unruhe

optische Halluzinationen

Personenverkennung

Wechselnde Symptomatik!!!!

delirante Verwirrung 

organische Symptome:

Tachykardie

Mydriasis (Pupillen weit)

trockene rote Haut

Krämpfe


z.B. Entzugsdelir bei Alkohol, Drogen

\subparagraph*{
Ma{\ss}nahmen}


NAW, ad Interne \"Uberwachungsstation (RR  instabil!)


RR-Kontrollen

Realitätsverlust

\paragraph{Mögliche Ursachen}

Endogene Psychosen (Schizophrenie,  schizoaffektive Psychose)

Organische Psychosen (Alkoholparanoia,  Drogenpsychose (z.B. Coke bugs
),  Encephalitis)

Realitätsverlust - Wahn

Wahn = falsche, allerdings in sich  logische Auffassung der Realität,
an der  unkorrigierbar festgehalten wird  (paranoid = wahnhaft) 

Themen: Verfolgungswahn, Grö{\ss}enwahn,  Verarmungswahn,...


Realitätsverlust - Halluzination


Halluzinationen: Trugwahrnehmungen,  f\"ur alle Sinne möglich
(wei{\ss}e Mäuse,  Stimmen,
{\quotedblbase}Geisterber\"uhrungen{\textquotedblleft},...)


Psychomotorische Unruhe /  Aggressivität


Organische Ursachen ausschlie{\ss}en!!!

{\mdseries\color[rgb]{0.0,0.0,0.5019608}
Intoxikation durch Alkohol, Drogen,  Medikamente}


Hypoglykämie

St.p. SHT

Hirntumor

psychomotorische Unruhe /  Aggressivität

Schwere Erregungszustände:

Katatonie : Starre mit Sprechstereotypien, Grimassieren,
Befehlsautomatismen, Patient imitiert Untersucher

Manie: Distanzlosigkeit, aggressiv gefärbte Agitiertheit

psychoreaktiv: abnorme Persönlichkeit

psychosoziale Krise: Reaktion auf Anla{\ss} / Stre{\ss}reaktion

psychomotorische Unruhe /  Aggressivität

Hysterie (funktionelle Erregungszustände):

anla{\ss}gebunden, entsprechende Grundpersönlichkeit

psychosomatischer Funktionsverlust (Lähmungen, Erblindung) ohne echte
organische Ursache

Entspricht psychischer Abwehrreaktion

{\mdseries\color[rgb]{0.0,0.0,0.5019608}
hysterischer Krampfanfall:}

{\mdseries
keine Verletzungen, keine Apnoe, atypische Krämpfe, kein
Stuhl-/Harnverlust}

{\mdseries
unbewu{\ss}t {\quotedblbase}gespielt{\textquotedblleft}}

{\mdseries
psychomotorische Unruhe /  Aggressivität}

\paragraph{Ma{\ss}nahmen:}

{\mdseries
ruhig bleiben! Achtung auf eigenen Tonfall!}

{\mdseries
NICHT provozieren!}

{\mdseries
falls Patient bewaffnet, ist Polizei zuständig}

{\mdseries
Selbstschutz besonders wichtig!!!}

{\mdseries
Manie}

{\mdseries
rastlose Unruhe - {\quotedblbase}energiegeladen{\textquotedblleft}}

{\mdseries
gereiztes Verhalten}

{\mdseries
\"Uberaktivität}

{\mdseries
Distanzverlust}

{\mdseries
Verwirrungszustände}

{\mdseries
Aggression}

{\mdseries
Selbst\"uberschätzung}

{\mdseries
ev. als bipolare Störung abwechselnd mit  Depression}

{\mdseries
Depression}

{\mdseries
Traurigkeit}

{\mdseries
Schuldgef\"uhle}

{\mdseries
fehlender Antrieb}

{\mdseries
Schlafstörungen}

{\mdseries
Suizidgedanken}

{\mdseries
Depression / Verzweiflung / Suizidalität}

{\mdseries
bei vielen psychiatrischen, aber auch  organischen Erkrankungen}

{\mdseries
Grundstimmung entweder:}

{\mdseries
antriebslos  oder}

{\mdseries
agitiert (höhere Suizid gefahr!)}

\paragraph{Ma{\ss}nahmen:}

{\mdseries
auf Gegenwart konzentrieren, soziale  Umgebung miteinbeziehen, 
Einf\"uhlungsvermögen!!!}

{\mdseries
Depression / Verzweiflung / Suizidalität}

{\mdseries
Selbstmordversuch ={\textgreater} stationäre  Aufnahme, falls
notwendig Parere!!!}

{\mdseries
Selbstmordphantasien: je konkreter,  desto gefährlicher!}

{\mdseries
Suizid = Selbstmord}

{\mdseries
Achtung, ev. sogar Mitnahmeselbstmord  möglich!}

{\mdseries
Angst}

{\mdseries
Gef\"uhl der Bedrohung mit vegetativer  Begleitsymptomatik}

{\mdseries
oft ohne erkennbaren Grund}

{\mdseries
bei psychiatrischen und organischen  Krankheitsbildern}

{\mdseries
Angst}



\subsection{Panikattacke}

\begin{itemize}
    \item plötzliches Angstgef\"uhl
    \item Dauer: Sekunden bis Minuten
    \item Thoraxbeschwerden (Infarktangst)
    \item Hyperventilation
    \item organisch mehrfach durchuntersuchte  Patienten, wirken gesund
    \item Erwartungsangst
    \item Vermeidungsverhalten
    \item Antriebslosigkeit / Stupor
\end{itemize}



\subsection{Stupor = Erschlaffung}

\paragraph{Symptome}

\begin{itemize}
    \item ausdrucksloser Patient
    \item unbeweglich
    \item starre Mimik
    \item spricht nicht, i{\ss}t nicht ={\textgreater} ev. bis Verhungern!
    \item Suchterkrankungen
    \item primär psychisches Problem
    \item sekundär körperliche und soziale Folgen
    \item eigengesetzlicher Ablauf
    \item zunehmender Kontrollverlust \"uber  eigenes Verhalten
\end{itemize}



\subsection{Alkohol- und Medikamentenentzug}


\paragraph{Symptome}

\begin{itemize}
    \item Schwitzen
    \item Unruhe
    \item Tremor  (Zittern der Hände)
    \item Delirium  (Bewu{\ss}tseinsstörungen, Halluzinationen )
    \item Krampfanfälle
    \item RR-Entgleisung
    \item {\quotedblbase}Cold turkey{\textquotedblleft} bei Heroinentzug
\end{itemize}


\paragraph{Anamnese}

\begin{itemize}
    \item Plötzlich kein Zugang mehr zu Alkohol:
    \item K\"urzlicher Einzug in Altersheim
    \item k\"urzlich pflegebed\"urftig geworden
    \item \ldots
\end{itemize}


\paragraph{Therapie}

\begin{itemize}
    \item Ruhige Gesprächsf\"uhrung, keine Kritik
    \item Suchtanamnese
\end{itemize}


\paragraph{Notarztindikation}

\begin{itemize}
    \item Krampfanfälle
    \item RR-Entgleisung
    \item Transportverweigerung
\end{itemize}



%%%%%%%%%%%%%%%%%%%%%%%%%%%%%%%%%%%%%%%%%%%%%%%%%%%%%%%%%%%

\section{Systemische Hitzeschäden}

Der menschliche Körper hält seine Kerntemperatur  im Normalfall
zwischen 36,5{\textdegree}  und 37,5{\textdegree} C  konstant.
\"Ubersteigt  die Wärmezufuhr die effektive  Wärmeabgabe,
resultiert ein Hitzeschaden: \gT{Hitzekollaps} , 
\gT{Hitzeerschöpfung} , \gT{Hitzschlag}
 (bzw. bei  direkter Sonneneinstrahlung auf den Kopf ein  Sonnen\-stich
).



\subsection{Hitzekollaps}

\gMarried{Bei einem Wärmestau, bei dem die zugeführte Wärme
größer ist als die abgegebene, kann es zu einem so
genannten Hitzekollaps kommen. Die Blutgefäße in der Haut erweitern sich und das Blut versackt, der Blutdruck fällt und die
Durchblutung des Gehirns nimmt kurzfristig ab. Die Folge ist eine kurze Ohnmacht,
den Patient wird "`schwarz vor Augen"'. Lagert man den Patienten flach oder
eventuell mit erhöhten Beinen, kühlt ihn beziehungsweise
bringt ihn an einen kühlen Ort, so bessert sich sein Zustand recht rasch. } 
{\begin{itemize} \item
\textbf{Wärmestau}  (zugef\"uhrte Wärme  kurzfristig {\textgreater} abgegebene)
\item Erweiterung der Blutgefä{\ss}e in der Haut  (maximal)
\item RR\,${\downarrow}$ , Durchblutung des Gehirns
kurzzeitig ${\downarrow}$
\item kurze Ohnmacht, {\quotedblbase}schwarz vor Augen{\textquotedblleft}
\end{itemize}}

\begin{itemize}
\item Flachlagerung: ${\rightarrow}$ Gehirndurchblutung
wieder ${\uparrow}$, bewu{\ss}tseinsklar
\item nach Sekunden gebessert, {\quotedblbase}regelt sich von
selbst{\textquotedblleft}
\end{itemize}
\begin{itemize}
\item ${\rightarrow}$ Therapie: Patient sitzen lassen,
Fl\"ussigkeit, ev. BZ messen
\end{itemize}



\subsection{Hitzeerschöpfung}

\gMarried{Bei der Hitzeerschöpfung ist die Wärmeaufnahme
dauerhaft größer als die Wärmeabgabe. Der Körper versucht
dies durch starkes schwitzen zu kompensieren, dadurch kommt 
es zu einem starken Wasser und Elektrolytverlust. Der
Blutdruck fällt, es kommt zu einer Tachykardie, die
Körpertemperatur bleibt jedoch noch normal. 
Hier ist es wichtig den Patienten zu
kühlen.}{\begin{itemize}
\item Wärmeaufnahme {\textgreater}  Wärmeabgabe
\item durch starkes Schwitzen Wasser- und  Elektrolytverlust
\item starkes Durstgef\"uhl
\item RR ${\downarrow}$, Tachykardie
\item Körpertemperatur normal
\end{itemize}}




\begin{gMassnahmenEnv}
\paragraph{Therapie}

\begin{itemize}
\item
K\"uhlung, Schatten,  Fl\"ussigkeitszufuhr, Vitalparameter,
evt. BZ messen
\end{itemize}
\end{gMassnahmenEnv}



\subsection{Hitzschlag}

\gMarried{Beim Hitzeschlag ist nun auch die Regulation der Körpertemperatur gestört, es
sind auch sehr hohe Werte (bis 43°) möglich. Dies bedeutet akuter Lebensgefahr.
Der Patient hat eine zunächst trockene, rote und weiße,
später grau-bläuliche Haut, der Kopf ist oft hochrot.
Weiters kommt es zu neurologischen Symptomen wie massiven
Kopfschmerzen, Schwindel, Erbrechen und
Bewusstseinstörungen; bis hin zur Bewusstlosigkeit. Der Patient 
kann Schockzeichnen zeigen, evt. auch eine schnelle, flache
Atmung.}{
\begin{itemize}
    \item Kopfschmerz, Schwindel, Erbrechen
    \item Haut zunächst trocken, rot und hei{\ss},  später grau-blau, Kopf
    hochrot
    \item stark  erhöhte  Körpertemperatur (bis  43{\textdegree}C möglich
    ={\textgreater} LEBENSGEFAHR !!!)
    \item Bewusstseinsstörung, Bewusstlosigkeit
    \item Schockzeichen
    \item schnelle, flache Atmung
\end{itemize}}

\begin{gMassnahmenEnv}
\paragraph{Ma{\ss}nahmen}

\begin{itemize}
\item Flachlagerung an k\"uhlem Ort (Schatten), Beine und Kopf hochlagern
\item beengende Kleidungsst\"ucke lockern
\item K\"uhlung von au{\ss}en (Fächer, mit Eisw\"urfeln abreiben etc.)
\item K\"uhlung von innen: f\"ur ANSPRECHBARE Patienten viel k\"uhle
(alkoholfreie) Fl\"ussigkeit, ev. Elektrolyt-Getränk
\item Schocklagerung u. -bekämpfung
\item Vitalparameter, $O_2$, BZ messen
\end{itemize}
\end{gMassnahmenEnv}



\subsection{Sonnenstich}

\gDefinition{Hitzeschaden durch Hirnerwärmung infolge von
direkter Sonneneinstrahlung auf den \textbf{ungesch\"utzten
Kopf}. }

\gMarried{Der Sonnenstich ist eine Sonderform des Hitzeschadens: 
er entsteht durch direkte Sonneneinstrahlung auf den
ungeschützten Kopf. Dadurch steigt die Temperatur im Kopf,
es kommt zu einer Reizung der Hirnhäute: 

Symptome wie bei
einer Hirnhautentzündung (Meningitis) beziehungsweise andere
neurologische Symptome wie zum Beispiel Kopfschmerz sind die
Folge. Die Gesichts- und Kopfhaut ist hochrot und heiß, der 
restliche Körper bleibt dagegen eher kühl. Der Patient
wirkt abgeschlagen und klagt über heftige 
Kopfschmerzen, Schwindel und Übelkeit, auch Unruhe, 
Verwirrtheitszustände und Nackensteifigkeit sind zu
beobachten. In schweren Fällen kann es sogar bis zu 
Krampfanfällen und Bewusstlosigkeit kommen.}
{\begin{itemize}
\item {\mdseries
direkte  Sonneneinstrahlung auf den \textbf{ungesch\"utzten 
Kopf}${\rightarrow}$ Temperatur  im Gehirn  steigt  ${\rightarrow}$
Reizung der Hirnhäute  ${\rightarrow}$ Symptome wie bei 
Hirnhautentz\"undung (Meningitis) bzw.  neurologische Symptome durch 
Schwellung des ganzen Gehirnes}
\item {\mdseries
Kann kombiniert mit anderen Hitzeschäden auftreten}
\item {\mdseries
Gesichts- u. Kopfhaut hochrot u. hei{\ss}, restlicher Körper k\"uhl}
\item {\mdseries
Abgeschlagenheit, heftiger Kopfschmerz, Schwindel}
\item {\mdseries
Unruhe,Verwirrtheitszustände, Brechreiz}
\item {\mdseries
Nackensteifigkeit}
\item {\mdseries
Schwere Fälle: Krampfanfälle, Bewusstlosigkeit}
\end{itemize}}




\begin{gMassnahmenEnv}
\paragraph{Ma{\ss}nahmen}

\begin{itemize}
    \item Flachlagerung an k\"uhlem Ort (Schatten), Kopf hochlagern
    \item beengende Kleidungsst\"ucke lockern
    \item K\"uhlung von au{\ss}en (Fächer, mit Eisw\"urfeln abreiben, kalte
    Umschläge auf die Stirn etc.)
    \item K\"uhlung von innen: f\"ur ANSPRECHBARE Patienten viel k\"uhle
    (alkoholfreie) Fl\"ussigkeit, ev. Elektrolyt-Getränk
    \item Schocklagerung u. -bekämpfung
    \item Vitalparameter, $O_2$, BZ messen
\end{itemize}
\end{gMassnahmenEnv}



\section{Unterk\"uhlung}

\"Ahnlich wie beim Schock kommt es auch  bei der Unterk\"uhlung zur
Zentralisation  des Blutes, um die  Körperkerntemperatur aufrecht
erhalten  zu können. Je niedriger diese  Kerntemperatur jedoch sinkt,
desto  schwerwiegender sind die Folgen!

\paragraph{Phasen}

\begin{itemize}
\item {\mdseries
\gT{Kampfphase}:  36,5{\textdegree}-34{\textdegree}C,
Kältezittern, Herzfrequenz steigt ({\textgreater}100
Schläge/Minute), schnelle Atmung, Schmerzen}

\item {\mdseries
\gT{Erschöpfungsphase}:
34{\textdegree}-30{\textdegree}C, Bewusstseinstr\"ubung, Kältestarre,
Herzfrequenz sinkt wieder ({\textless}60/min), zu langsame Atmung,
keine Schmerzen mehr}

\item {\mdseries
\gT{Kältenarkose}: 30{\textdegree}-27{\textdegree}C,
Bewusstlosigkeit, unregelmä{\ss}iger Herzschlag, Atempausen}

\item {\mdseries
{\textless}27{\textdegree}C ${\rightarrow}$ klinisch tot}

\end{itemize}


\begin{gMassnahmenEnv}
\paragraph{Ma{\ss}nahmen}

\begin{itemize}
    \item nur soviel wie UNBEDINGT nötig ist  bewegen (BERGUNGSTOD)
    \item Schutz vor weiterer Abk\"uhlung, KEINE AKTIVE ERW\"ARMUNG!
    \item Woll- + Aludecken! M\"utze!
    \item ev. Wiederbelebungsma{\ss}nahmen
    \item Notarzt !
\end{itemize}
\end{gMassnahmenEnv}



Bei Unterk\"uhlung gilt:

\gFazit{{\quotedblbase}Nobody is dead until warm and dead.{\textquotedblleft}

Schwer unterk\"uhlte Patienten können auch noch nach Stunden (!) ohne
Kreislauftätigkeit erfolgreich reanimiert werden!}



%%%%%%%%%%%%%%%%%%%%%%%%%%%%%%%%%%%%%%%%%%%%%%%%%%%%%%%%%%%
\section{Vergiftungen}

Michael Withofner, Dr. med. univ.

\subparagraph*{
Vergiftung - Definition}

{\mdseries
\gT{Vergiftung} bedeutet eine (meist  dosisabhängige)
organische  Funktionsbeeinträchtigung / Schädigung  des Organismus
durch eine (organische/ anorganische) chemische Substanz!}

{\mdseries
(z.B.: NaCl ={\textgreater} 20 dag tödlich,  Medikamente
={\textgreater} teilweise im Mikrogramm- Bereich tödlich!!!)}

{\mdseries
Vergiftungen - erkennen}

(Au{\ss}en-)Anamnese besonders wichtig (Vergiftungsart/Ursache), oft
einziger Aufschlu{\ss}, da Patient ev. bewu{\ss}tseinsgetr\"ubt

\subparagraph*{6 W -Fragen:}

\begin{enumerate}
\item \gEs{Wer}? (Alter, KG, Vorerkrankungen,...
={\textgreater} Risikoabschätzung!)
\item \gEs{Was}? (Art des Giftes ={\textgreater} ev.
zu erwartende Symptome)
\item \gEs{Wann}? (Zeit seit Einnahme ={\textgreater}
Magensp\"ulung ja/nein?)
\item \gEs{Wie}?
(inhalativ/oral/perkutan/intravenös etc.)
\item \gEs{Wieviel}? (Dosisabschätzung)
\item \gEs{Warum}? (akzidentiell/absichtlich)
\end{enumerate}



\paragraph{Ursachen:}

{\mdseries
akzidentiell: z.B.}

(Arbeits-)Unfall: Reizgase, Chemikalien

Kinder: {\quotedblbase}Zuckerl{\textquotedblleft} = Medikamente

fahrlässig: falsch beschriftete Flaschen  (Putzmittel in
Getränkeflasche etc.)

{\mdseries
absichtlich:}

Suizid: meist Medikamente (Benzodiazepine,  Barbiturate)

Drogen: \"Uberdosis (gewollt oder ungewollt), 
{\quotedblbase}Verschnitt{\textquotedblleft}

Mord: kriminelle Absicht, bei Verdacht Polizei  einschalten!!!

\subparagraph*{Aufnahme}

{\mdseries
oral : Verschlucken  der giftigen  Substanz }

{\mdseries
inhalativ : Einatmen}

{\mdseries
rauchen/schnupfen: Kombination oral /  inhalativ}

{\mdseries
transkutan : fettlöslich e Substanz kann  Hautbarriere durchdringen
(z.B.:  organische Lösungsmittel,  Insektenschutzmittel etc.)}

{\mdseries
intravenös/intramuskulär/subkutan: spritzen}

{\mdseries
Vergiftungen -- unspezifische  Symptome}

{\mdseries
gastrointestinal:}

{\mdseries
Erbrechen}

{\mdseries
Durchfall (Diarrhoe )}

{\mdseries
Schmerzen im GIT}

{\mdseries
ZNS:}

{\mdseries
Bewu{\ss}tsein [F0EA?][F020?](Rausch)}

{\mdseries
ev. Koma}

{\mdseries
zerebrale Krampfanfälle}

\subparagraph*{Vorgehensweise}

{\mdseries
5-Fingerregel:}

\begin{enumerate}
\item Vitalparameter
\item Giftentfernung
\item Antidot-Therapie (Notarzt)
\item Gift-Asservierung
\item Transport
\end{enumerate}


\begin{gMassnahmenEnv}
\subparagraph*{Vergiftungen - Basisma{\ss}nahmen}

{\mdseries
allgemein:}

{\mdseries
Vitalparameter ständig \"uberwachen}

{\mdseries
O2  !!!}

{\mdseries
Asservierung (Gift bekannt -- ev. spezifische  Ma{\ss}nahmen!)}

{\mdseries
was liegt herum?}

{\mdseries
Mistk\"ubel durchsuchen}

{\mdseries
Erbrochenes (in Nierentasse mitnehmen!)}

{\mdseries
Vergiftungen -- Erstma{\ss}nahmen}

{\mdseries
primäre Giftentfernung / stoppen der Giftaufnahme}

{\mdseries
oral : }

{\mdseries
Patient soll nichts essen / trinken}

{\mdseries
NICHT absichtlich Erbrechen herbeif\"uhren}

{\mdseries
im Spital: Aktivkohle (häufig) bzw.  Magensp\"ulung (selten)}

{\mdseries
inhalativ :}

{\mdseries
Eigenschutz!!!}

{\mdseries
Entfernung aus Giftatmosphäre}

{\mdseries
O2  / Frischluft}

{\mdseries
Vergiftungen -- Erstma{\ss}nahmen}

{\mdseries
primäre Giftentfernung / stoppen der Giftaufnahme}

{\mdseries
transkutan :}

{\mdseries
Eigenschutz!!!}

{\mdseries
Entfernung kontaminierter Kleidung}

{\mdseries
Sp\"ulung der betroffenen Hautpartien mit Wasser  (richtig abrinnen
lassen)}

{\mdseries
okulär  (Auge):}

{\mdseries
Sp\"ulung mit Wasser / Ringerlösung (nach au{\ss}en abrinnen lassen!)}

{\mdseries
Kontaktlinsen entfernen}

{\mdseries
Vergiftungen -- spezifische  Ma{\ss}nahmen}

{\mdseries
bei bekanntem Gift ev. Antidot  (Gegengift)-Gabe möglich (Notarzt / 
Spital)}
\end{gMassnahmenEnv}

\subparagraph*{
Vergiftungszentrale:}

{\mdseries
406 43 43}

{\mdseries
bei bekanntem Giftstoff (6 W-Fragen) ev. spezifische Anweisungen
telefonisch}

{\mdseries
spezielle Vergiftungen}

{\mdseries
eventuell anhand von Symptomen  diagnostizierbar}

{\mdseries
{\quotedblbase}Toxidrome{\textquotedblleft} ={\textgreater} allerdings
eigene  Wissenschaft}

{\mdseries
dennoch gut, einige {\quotedblbase}typische{\textquotedblleft} 
Vergiftungen zu kennen}

{\mdseries
Medikamente / Drogen}

{\mdseries
Benzodiazepine  (Tranquilizer,  Schlafmittel):}

{\mdseries
Substanzen: Valium{\textregistered} , Rohypnol{\textregistered} ,
Dormicum{\textregistered} ,  Temesta{\textregistered}  ...}

{\mdseries
Symptome:}

{\mdseries
Bewu{\ss}tseinseintr\"ubung}

{\mdseries
RR [F0EA?]}

{\mdseries
Spontanatmung mä{\ss}ig [F0EA?]}

{\mdseries
Antidot: Anexate{\textregistered} }

{\mdseries
Medikamente / Drogen}

{\mdseries
Schmerzmittel (NSAR, Nicht Steroidale AntiRheumatika)}

{\mdseries
Substanzen: }

{\mdseries
Paracetamol: Mexalen{\textregistered}}

{\mdseries
Salicylate: Aspirin{\textregistered} , Brufen{\textregistered} ,
Voltaren{\textregistered}}

{\mdseries
Symptome / Folgen:}

{\mdseries
Paracetamol: }

{\mdseries
Leberschäden}

{\mdseries
Salicylate: }

{\mdseries
Azidose}

{\mdseries
Hyperthermie}

{\mdseries
Krämpfe}

{\mdseries
Therapie: Antidote im AKH / WSP lagernd,  Hämofiltration!!!}

{\mdseries
Medikamente / Drogen}

{\mdseries
Opiate  (aus Milchsaft des Schlafmohns (Papaver somniferum))}

{\mdseries
Substanzen: Morphin (Fentanyl{\textregistered} etc.), Codein, Heroin
(illegal)}

{\mdseries
Unterliegen strengem Suchtgiftgesetz (Schmerzmittel)}

{\mdseries
Schwarzmarkt (Heroin), {\quotedblbase}Fixer{\textquotedblleft}, sehr
hohes Suchtpotential}

\subparagraph*{Symptome:}

{\mdseries
Miose (Pupillen eng, stecknadelkopfgro{\ss}, nicht  lichtreaktiv)}

{\mdseries
Bewu{\ss}tsein [F0EA?]  bis Bewu{\ss}tlosigkeit}

{\mdseries
Atemdepression bis Atemstillstand (Brady- /Apnoe )}

{\mdseries
Therapie:}

{\mdseries
O2 , Beatmung, Vorgehen bei Bewu{\ss}tlosigkeit}

{\mdseries
Antidot: Narcanti{\textregistered}  (Wirkstoff: Naloxon )}

{\mdseries
Medikamente / Drogen}

{\mdseries
Achtung bei Kontakt mit drogenkranken  Patienten:}

{\mdseries
Immer  an Infektionserkrankung denken!!!}

{\mdseries
TBC}

{\mdseries
Hepatitis B + C}

{\mdseries
HIV}

{\mdseries
Immer Handschuhe tragen!!!}

{\mdseries
Blutkontakt unbedingt vermeiden}

{\mdseries
Händedesinfektion }

{\mdseries
Medikamente / Drogen}

\subsubsection{Intoxikation mit Alkohol}
{\mdseries
Substanzen: Methanol (Fuselalkohol,
{\quotedblbase}Vorlauf{\textquotedblleft}), Ethanol
({\quotedblbase}Alk{\textquotedblleft})}

{\mdseries
SEHR häufige Vergiftung (gesellschaftlich akzeptiert),
{\quotedblbase}Kampfsaufen{\textquotedblleft}!}

\subparagraph*{Symptome: }

{\mdseries
Promilleabhängig}

{\mdseries
Euphorie}

{\mdseries
Sprach-/ Gangunsicherheit, Enthemmung}

{\mdseries
Bewu{\ss}tseinstr\"ubung,  Schmerzunempfindlichkeit, 
Temperaturregulationsstörung, Inkontinenz}

{\mdseries
Bewu{\ss}tlosigkeit, Atemdepression, Aspiration,  Unterk\"uhlung}

\subparagraph*{Therapie:}

{\mdseries
Vitalparameter}

{\mdseries
stabile Seitenlage / Reanimation}

{\mdseries
BZ messen (Kinder!)}

{\mdseries
vorsichtige Beatmung (Alkoholisierte erbrechen  leichter!)}

{\mdseries
Alkohol -- chronische Schäden}

{\mdseries
Leberschäden}

{\mdseries
Zirrhose}

{\mdseries
Aszites}

{\mdseries
Krampfadern in der unteren Speiseröhre  (Blutungsgefahr)}

{\mdseries
Blutgerinnung [F0EA?]}

{\mdseries
Vitaminmangel ={\textgreater} }

{\mdseries
intellektueller Abbau (Vit B)}

{\mdseries
Anämie (Vit B12 , Folsäure)}

{\mdseries
Herzrhythmusstörungen}

{\mdseries
Pankreatitis (Entz\"undung der  Bauchspeicheldr\"use)}

{\mdseries
Alkohol - Entzug}

\subparagraph{Entzugsdelir}
{\mdseries
{\quotedblbase}Delirium tremens {\textquotedblleft}}

{\mdseries
Zittern (Tremor )}

{\mdseries
Unruhe, Angst, Psychose}

{\mdseries
Halluzinationen ({\quotedblbase}wei{\ss}e
Mäuse{\textquotedblleft},...)}

{\mdseries
Tachykardie}

{\mdseries
RR-Entgleisung}

{\mdseries
zerebrale Krampfanfälle}

{\mdseries
lebensbedrohlicher Zustand}

{\mdseries
{\quotedblbase}Akut-{\textquotedblleft}Therapie: Alkohol (Ethanol als
Schnaps  etc.)}

{\mdseries
im Spital: medikamentös unterst\"utzter Entzug }

{\mdseries\color[rgb]{0.0,0.0,0.5019608}
Fl\"ussigkeiten}

\subsubsection{Korrosiva (Säuren/Laugen):}
{\mdseries
Verätzung}

{\mdseries
Lokale Schädigung durch irreversible Zerstörung (Denaturierung) von
Eiwei{\ss}stoffen. }

{\mdseries
betroffenen Organe}

{\mdseries
GIT}

{\mdseries
Haut}

{\mdseries
Augen}

{\mdseries
Säuren/Laugen}

\subparagraph{Symptome bei Ingestion:}
{\mdseries
\"Atzspuren (Mund/Rachen), Abrinnspuren}

{\mdseries
W\"urgen}

{\mdseries
Erbrechen}

{\mdseries
Schmerzen}

{\mdseries
Vorsicht: Ruptur/Perforation/Blutungsgefahr!!!}

\begin{gMassnahmenEnv}
\subparagraph*{Therapie:}

\begin{itemize}
\item schluckweise Wasser trinken lassen
\item KEINE Neutralisationsversuche!!! (Milch  obsolet!!!)
\item NIEMALS absichtlich Erbrechen  herbeif\"uhren!!!!
\end{itemize}
\end{gMassnahmenEnv}


{\mdseries
Säuren/Laugen}

\subsubsection{Verätzung der Haut}
{\mdseries
Schorfbildung}

{\mdseries
glasige Verquellung}

{\mdseries
Nekrosen}

{\mdseries
Schmerzen}

\subparagraph*{
Therapie:}

{\mdseries
Eigenschutz! (Handschuhe)}

{\mdseries
kontaminierte Kleidung entfernen}

{\mdseries
betroffene Hautpartien mit Wasser absp\"ulen  (unverletzte Haut
sch\"utzen!)}

{\mdseries
locker, steril verbinden}

{\mdseries
Säuren/Laugen}

\subsubsection{Verätzung der Augen}
{\mdseries
Hornhauttr\"ubung, später Narben}

{\mdseries
Visusverlust}

{\mdseries
Schmerzen}



\begin{gMassnahmenEnv}
\subparagraph*{
Therapie:}

\begin{itemize}
\item {\mdseries\color[rgb]{0.0,0.0,0.5019608}
Sofortige Sp\"ulung mit Wasser / Ringerlösung}

\item {\mdseries\color[rgb]{0.0,0.0,0.5019608}
danach beide Augen keimfrei abdecken!}

\item {\mdseries\color[rgb]{0.0,0.0,0.5019608}
Fl\"ussigkeiten}

\end{itemize}
\end{gMassnahmenEnv}


\paragraph[Schaumbildner (Wasch-/Putzmittel)]{Schaumbildner
(Wasch-/Putzmittel)}
\subparagraph*{
Symptome:}

{\mdseries
verminderte Oberflächenspannung:  Aspirationsgefahr ${\uparrow}$}

{\mdseries
Verlegung der Atemwege}

{\mdseries
\"Ubelkeit/Erbrechen}

\subparagraph*{
Therapie:}

{\mdseries
Patienten n\"uchtern lassen}

{\mdseries
NIEMALS absichtlich Erbrechen  herbeif\"uhren!!!!}





\subsection{Stickgase}

\subparagraph*{
Bei Gasunfällen hat der \gEs{Selbstschutz} Vorrang.
Oft ist die Bergung nur mit \gEs{schwerem
Atemschutz} möglich ${\rightarrow}$ FEUERWEHR!}

{\mdseries
Ein toter Helfer ist kein Helfer.}

\paragraph{Kohlenmonoxid (CO)}
entsteht bei unvollständiger Verbrennung (O2- Zufuhr ${\downarrow}$)



{\mdseries
300x höhere Affinität zu Hämoglobin (als O2), }

{\mdseries
Sättigung beim Pulsoxymeter falsch hoch!!!}

\subparagraph*{
Symptome:}

{\mdseries
Dyspnoe  ohne  Zyanose  (rosige Hautfarbe, kirschrote Lippen!!!)}

{\mdseries
Tachykardie ={\textgreater} Gefahr: MCI}

{\mdseries
Schwindel,Kopfschmerz, \"Ubelkeit, Erbrechen}

{\mdseries
ev. Krämpfe bis Koma}

\subparagraph*{Therapie:}

{\mdseries
Entfernung aus Gasatmosphäre  (Eigenschutz!!!)}

{\mdseries
O2 hochdosiert , Hyperventilation, hyperbare  O2 -Therapie}

\begin{figure}[H]
    \begin{center}
\includegraphics[width=\linewidth]{./images/2006-01-31-gaswienheute2.png}
\captionof{figure}{\gAuthNotesPicLegalOk{Gas Wien Heute
Homepage}{OK, wenn richtig titiert}$CO$-Vergiftung.
Regelmäßig kommt es zu Unfällen durch defekte Gasthermen
oder Heizlüfter. Bei diesem Einsatz war ein RTW des ASB
Floridsdorf-Doanustadt beteiligt. Die dritte Patientin
wurde erst nachdem die Feuerwehr \gE{sämtliche} (!)
Wohnungen des Wohnhauses aufgebrochen hatte in der
\gE{darüberliegenden Wohnung} gefunden.
\gSourcePicture{Wien heute \url{http://wien.orf.at}}}
\end{center}
\end{figure}

\paragraph[Kohlendioxid (CO2) ]{Kohlendioxid (CO\textsubscript{2}) }
{\mdseries
Gärgas, Verbrennungsgas, entsteht auch bei  Atmung}

{\mdseries
schwerer als Luft (sammelt sich in Mulden,  Weinkellern etc.)}

{\mdseries
Symptome:}

{\mdseries
Kopfschmerz, Schwindel, \"Ubelkeit, Krämpfe, Ersticken}

{\mdseries
O2-Aufnahme ${\downarrow}$, stimuliert Atemzentrum (Hyperventilation)}

\subparagraph*{
Therapie}

{\mdseries
Bergung (Eigenschutz), Frischluft}

{\mdseries
10-15 l O2  }

{\mdseries\color[rgb]{0.0,0.0,0.5019608}
Reizgase}

{\mdseries
Soforttyp : (z.B. Tränengas,...)}

{\mdseries
Konjunktivitis}

{\mdseries
starker Hustenreiz}

{\mdseries
Verzögerungstyp: (z.B. Nitrosegase)}

{\mdseries
kaum Hustenreiz}

{\mdseries
Latenzphase (kaum Symptome)}

{\mdseries
3-24h später: toxisches Lungenödem}

{\mdseries
bei Nichtbehandlung ev. irreversible Lungenfibrose}

{\mdseries
Reizgase}

\paragraph{Vorkommen}

{\mdseries
Chemische Industrie: Nitrosegase,  Schwefelwasserstoffe, Ammoniak}

{\mdseries
 Kunststoffindustrie: Chlorwasserstoffe,  Formaldehyd}

{\mdseries
D\"ungemittelherstellung: Ammoniak,  Nitrosegase, Schwefelwaserstoffe}

{\mdseries
Haushalt: Chlorgas (Kombination von best.  Haushaltsreinigern mit
Essig), Lacke,  Imprägniersprays, ...}

{\mdseries
Reizgase}

{\mdseries
Therapie:}

{\mdseries
Bergung (Eigenschutz!)}

{\mdseries
100\% O2 , ev. PEEP-Beatmung durch NFS /  Notarzt}

{\mdseries
wiederholt inhalatives Corticoid (Pulmicort{\textregistered} )  durch
NFS / Notarzt}

{\mdseries
Bei Reizgas / Stickgas immer auch nach  anderen bewu{\ss}tlosen Personen
fragen /  suchen!!! (unter Eigenschutz!!!)}

\putbib



\clearpage
